\documentclass[italian,12pt]{article} %tipo di documento

\def\Title{Verbale interno}
\def\Subtitle{Riunione interna settimanale}
\def\Author{7Last}
\def\Date{2024-05-22}
\def\Version{v1.0}

%pacchetti extra da scaricare dblfloatfix, fancyhdr
\usepackage[left=2cm, right=2cm, bottom=3cm, top=3cm]{geometry}
\usepackage{fancyhdr}%creazione header-footer
\usepackage{graphicx} %serve per inserire immagini
\graphicspath{ {../../../logo/} }
%\usepackage{dblfloatfix} %serve per posizionare gli elementi dove si vuole
\usepackage[hidelinks]{hyperref} %serve per i link
\usepackage{tikz}
\usepackage{tgadventor} % font
\usepackage[useregional=numeric,showseconds=true,showzone=false]{datetime2}
\usepackage{caption}

\usepackage{hyperref}
\usepackage{tocloft}
\usepackage{titlesec}
\usepackage{color}
\usepackage{ulem}
\usepackage{pgfplots}
\usepackage{pgf-pie}
\usepackage[italian]{babel}
\usepackage{comment}
\usepackage{tabularx}
\usepackage{longtable}
\usepackage{float}
\usepackage{amsmath}
% Definizione delle nuove classi di titolo
\titleclass{\subsubsubsection}{straight}[\subsection]
\titleclass{\subsubsubsubsection}{straight}[\subsubsubsection]
\titleclass{\subsubsubsubsubsection}{straight}[\subsubsubsubsection] % nuovo livello

% Creazione dei nuovi contatori
\newcounter{subsubsubsection}[subsubsection]
\newcounter{subsubsubsubsection}[subsubsubsection]
\newcounter{subsubsubsubsubsection}[subsubsubsubsection] % nuovo livello

% Rinnovo dei comandi per la formattazione dei numeri delle sezioni
\renewcommand\thesubsubsubsection{\thesubsubsection.\arabic{subsubsubsection}}
\renewcommand\thesubsubsubsubsection{\thesubsubsubsection.\arabic{subsubsubsubsection}}
\renewcommand\thesubsubsubsubsubsection{\thesubsubsubsubsection.\arabic{subsubsubsubsubsection}} % nuovo livello
\renewcommand\theparagraph{\thesubsubsubsubsubsection.\arabic{paragraph}} % opzionale; utile se i paragrafi devono essere numerati

% Formattazione dei titoli delle sezioni
\titleformat{\subsubsubsection}
  {\normalfont\normalsize\bfseries}{\thesubsubsubsection}{1em}{}
\titleformat{\subsubsubsubsection}
  {\normalfont\normalsize\bfseries}{\thesubsubsubsubsection}{1em}{}
\titleformat{\subsubsubsubsubsection} % nuovo livello
  {\normalfont\normalsize\bfseries}{\thesubsubsubsubsubsection}{1em}{} 

% Spaziatura dei titoli delle sezioni
\titlespacing*{\subsubsubsection}
{0pt}{3.25ex plus 1ex minus .2ex}{1.5ex plus .2ex}
\titlespacing*{\subsubsubsubsection}
{0pt}{3.25ex plus 1ex minus .2ex}{1.5ex plus .2ex}
\titlespacing*{\subsubsubsubsubsection} % nuovo livello
{0pt}{3.25ex plus 1ex minus .2ex}{1.5ex plus .2ex}

\makeatletter
% Rinnovo dei comandi per la formattazione dei paragrafi e sottoparagrafi
\renewcommand\paragraph{\@startsection{paragraph}{6}{\z@}%
  {3.25ex \@plus1ex \@minus.2ex}%
  {-1em}%
  {\normalfont\normalsize\bfseries}}
\renewcommand\subparagraph{\@startsection{subparagraph}{7}{\parindent}%
  {3.25ex \@plus1ex \@minus .2ex}%
  {-1em}%
  {\normalfont\normalsize\bfseries}}

% Definizione dei livelli per il Table of Contents
\def\toclevel@subsubsubsection{4}
\def\toclevel@subsubsubsubsection{5}
\def\toclevel@subsubsubsubsubsection{6} % nuovo livello
\def\toclevel@paragraph{7}
\def\toclevel@subparagraph{8}

% Definizione della formattazione per il Table of Contents
\def\l@subsubsubsection{\@dottedtocline{4}{7em}{4em}}
\def\l@subsubsubsubsection{\@dottedtocline{5}{10em}{5em}}
\def\l@subsubsubsubsubsection{\@dottedtocline{6}{14em}{6em}} % nuovo livello
\def\l@paragraph{\@dottedtocline{7}{18em}{7em}}
\def\l@subparagraph{\@dottedtocline{8}{22em}{8em}}
\makeatother

% Impostazione della profondità dei numeri di sezione e del Table of Contents
\setcounter{secnumdepth}{6} % nuovo livello
\setcounter{tocdepth}{6} % nuovo livello


\linespread{1.2}
\captionsetup[table]{name=Tabella}
\geometry{headsep=1.5cm}

\renewcommand{\contentsname}{Indice}%imposto il nome dell'indice
\renewcommand\familydefault{\sfdefault}

\renewcommand{\listtablename}{Indice delle tabelle}%imposto il nome della lista tabelle
\renewcommand\familydefault{\sfdefault}

\renewcommand{\listfigurename}{Indice delle immagini}%imposto il nome della lista immagini
\renewcommand\familydefault{\sfdefault}

\addto\captionsitalian{\renewcommand{\figurename}{Figura}}

%-------------------INIZIO DOCUMENTO--------------
\begin{document}

\newgeometry{left=2cm,right=2cm,bottom=2.1cm,top=2.1cm}
\begin{titlepage}
    \vspace*{.5cm}

    \vspace{2cm}
    {
        \centering
        {\bfseries\huge \Title\par}
        \bigbreak
        {\bfseries\large \Author\par}
        \bigbreak
        {\Version\par}
        \vfill

        \begin{tikzpicture}[remember picture,overlay]

            \fill[blue!20!black] (current page.south west) -- ++(10cm,0) -- ++(-10cm,15cm) -- cycle;
            \fill[orange] (current page.south east) -- ++(-18cm,0) -- ++(21.6cm,22cm) -- cycle;

            \clip (0,-3cm) circle (2.5cm) node (current page.center) {\includegraphics[width=5cm]{logo.jpg}};
        \end{tikzpicture}

    }

    \vfill

\end{titlepage}

\restoregeometry






















\newpage

%---------------header------------------%
\pagestyle{fancy}
\fancyhead{} % pulizia degli header
\lhead{%
\begin{tikzpicture}
    \clip (0,0) circle (0.5cm);
    \node at (0,0) {\includegraphics[width=1cm]{logo}};
\end{tikzpicture}%
}
\chead{\vspace{\fill}\Title\vspace{\fill}}
\rhead{\vspace{\fill}\Version\vspace{\fill}}
%quad è una spaziatura
%---------------------------------------%

%-----------tabella revisioni-----------%
\begin{table}[!h]
	\caption*{Versioni}
	\begin{center}
		\begin{tabular}{ l l l l l}
			\hline\\[-2ex]
			Ver. & Data       & Autore           & Verificatore   & Descrizione\\
			\\[-2ex] \hline \\[-1.5ex] \\
			0.3  & 2024-04-05 & Raul Seganfreddo &  & Aggiunta Introduzione		 \\
			0.2  & 2024-04-05 & Matteo Tiozzo    &  & Modificato tabella versioni\\
			0.1  & 2024-03-29 & Raul Seganfreddo &  & Creazione struttura documento\\
			\\[-1.5ex] \hline
		\end{tabular}
	\end{center}
\end{table}
%---------------------------------------%

\newpage

\tableofcontents

\listoftables

\listoffigures

\newpage

\section{Introduzione}
\subsection{Obiettivo del documento}
Il presente documento ha lo scopo di definire le strategie di verifica e validazione utilizzate per assicurare il corretto funzionamento dello strumento sviluppato e delle
attività che lo accompagnano.  Sarà sottoposto a revisioni continue, così da prevedere situazioni precedentemente non occorse e da seguire l'evoluzione del progetto.
\subsection{Glossario}
Il \href{https://7last.github.io/docs/rtb/documentazione-interna/glossario#glossario}{glossario\textsubscript{G}} è uno strumento utilizzato per risolvere eventuali dubbi riguardanti 
alcuni termini specifici utilizzati nella redazione del documento.
Esso conterrà la definizione dei termini evidenziati e sarà consultabile al seguente \href{https://7last.github.io/docs/rtb/documentazione-interna/glossario}{link}. I termini presenti in tale documento saranno evidenziati da una 'G' a pedice.
\subsection{Riferimenti}
\subsubsection{Riferimenti normativi}
\begin{itemize}
    \item \href{https://7last.github.io/docs/rtb/documentazione-interna/glossario#norme-di-progetto}{Norme di progetto\textsubscript{G}} (aggiungere versione e/o link al documento);
    \item Regolamento del progetto:\\
		  \url{https://www.math.unipd.it/~tullio/IS-1/2023/Dispense/PD2.pdf}.
\end{itemize}
\subsubsection{Riferimenti informativi}
\begin{itemize}
    \item \href{https://7last.github.io/docs/rtb/documentazione-interna/glossario#capitolato}{Capitolato\textsubscript{G}} d'appalto C6: \href{https://7last.github.io/docs/rtb/documentazione-interna/glossario#synccity}{SyncCity\textsubscript{G}} – A \href{https://7last.github.io/docs/rtb/documentazione-interna/glossario#smart-city}{smart city\textsubscript{G}} monitoring platform\\
    \url{https://www.math.unipd.it/~tullio/IS-1/2023/Progetto/C6.pdf};
    \item \href{https://it.wikipedia.org/wiki/ISO/IEC_9126}{Standard ISO/IEC 9126};
    \item \href{https://iso25000.com/index.php/en/iso-25000-standards/iso-25010}{Standard ISO/IEC 25010};
    \item \href{ https://en.wikipedia.org/wiki/ISO/IEC_12207}{Standard ISO/IEC 12207:1995};
    \item \href{URL}{\textit{Verbali esterni}};
    \item \href{URL}{\textit{Verbali interni}};
    \item \href{URL}{\href{https://7last.github.io/docs/rtb/documentazione-interna/glossario#analisi-dei-requisiti}{\textit{Analisi dei requisiti}\textsubscript{G}}};
    \item AGGIUNGERE LINK
\end{itemize}

\newpage

\section{Processi primari}
\subsection{Fornitura}
\subsubsection{Introduzione}
Il processo di fornitura è un percorso strutturato che stabilisce un accordo contrattuale tra il fornitore e il cliente e guida lo sviluppo e la consegna di un prodotto software, dalla concezione iniziale fino alla manutenzione post-rilascio. Questo processo è fondamentale per garantire che il software soddisfi i requisiti del cliente, sia di alta qualità e venga consegnato nei tempi e nei costi previsti.
\subsubsection{Attività}
Il processo di fornitura comprende le fasi di seguito specificate.
\begin{itemize}
	\item \textbf{Preparazione della proposta}: questa fase iniziale prevede la raccolta di informazioni e la stesura di una proposta formale per il cliente. Si suddivide in:
	      \begin{itemize}
		      \item analisi delle esigenze del cliente;
		      \item studio di fattibilità;
		      \item elaborazione della proposta.
	      \end{itemize}
	\item \textbf{Contrattazione}: durante questa fase il fornitore e il cliente discutono e negoziano i termini del contratto. Alcune attività tipiche di questa fase includono:
	      \begin{itemize}
		      \item discussione dei termini e delle condizioni;
		      \item stesura e revisione del contratto;
		      \item firma del contratto.
	      \end{itemize}
	\item \textbf{Pianificazione}: fondamentale per organizzare e programmare le attività del progetto. Questa fase include:
	      \begin{itemize}
		      \item stesura delle \href{https://7last.github.io/docs/pb/documentazione-interna/glossario\#milestone}{milestone\textsubscript{G}};
		      \item stesura del \href{https://7last.github.io/docs/pb/documentazione-interna/glossario\#piano-di-progetto}{piano di progetto\textsubscript{G}};
		      \item assegnazione dei compiti e delle risorse.
	      \end{itemize}
	\item \textbf{Esecuzione}: si effettua la realizzazione concreta del progetto, con la costruzione del prodotto software. L'attività è costituita principalmente da:
	      \begin{itemize}
		      \item sviluppo del software;
		      \item test e verifica;
		      \item documentazione.
	      \end{itemize}
	\item \textbf{Revisione}: prevede una valutazione approfondita del lavoro svolto per garantire che tutto sia conforme agli standard di qualità e ai requisiti contrattuali. Questa fase include:
	      \begin{itemize}
		      \item revisione del codice;
		      \item test di accettazione;
		      \item correzione delle discrepanze.
	      \end{itemize}
	\item \textbf{Consegna}: consiste, appunto, nella consegna del prodotto finale al cliente e nella preparazione per il supporto post-rilascio. Si compone di:
	      \begin{itemize}
		      \item consegna del software;
		      \item formazione del personale;
		      \item supporto post-rilascio.
	      \end{itemize}
\end{itemize}

\subsubsection{Contatti con l'azienda proponente}
\textit{Sync Lab S.r.l.} offre un indirizzo email e un canale Discord per la comunicazione tramite messaggi, oltre a Google Meet per gli incontri telematici. Gli incontri online si terranno inizialmente con cadenza bisettimanale, con la possibilità di organizzare incontri extra su richiesta del gruppo. Per ogni colloquio con l’azienda \href{https://7last.github.io/docs/pb/documentazione-interna/glossario\#proponente}{proponente\textsubscript{G}} sarà redatto un verbale che riassumerà i temi trattati. Tali documenti saranno disponibili al seguente \href{https://7last.github.io/docs/category/verbali-esterni}{\underline{link}}.

\subsubsection{Documentazione fornita}
Di seguito saranno elencati i documenti che il gruppo \textit{7last} consegnerà all'azienda \textit{Sync Lab S.r.l.} e ai \href{https://7last.github.io/docs/pb/documentazione-interna/glossario\#committente}{committenti\textsubscript{G}} \textit{Prof. Tullio Vardanega} e \textit{Prof. Riccardo Cardin}.
\newpage

\subsubsubsection{Piano di qualifica}
Il \href{https://7last.github.io/docs/pb/documentazione-interna/glossario\#piano-di-qualifica}{\textit{Piano di Qualifica}\textsubscript{G}} è un documento che specifica le responsabilità e le attività del \href{https://7last.github.io/docs/pb/documentazione-interna/glossario\#verificatore}{verificatore\textsubscript{G}} nel progetto, delineando le strategie e gli approcci adottati per garantire la qualità del prodotto software in fase di sviluppo. Redatto dal \href{https://7last.github.io/docs/pb/documentazione-interna/glossario\#verificatore}{verificatore\textsubscript{G}}, questo documento descrive le modalità di verifica e validazione, nonché gli standard e le procedure di qualità da seguire durante l'intero ciclo di vita del progetto. Tutti i membri del team si riferiranno a questo documento per assicurarsi di raggiungere la qualità desiderata. Questo documento è organizzato in diverse sezioni, tra cui:
\begin{itemize}
	\item \textbf{qualità di processo}: definisce gli standard e le procedure adottate durante il progetto al fine di garantire che i processi seguiti siano di alta qualità;
	\item \textbf{qualità di prodotto}: stabilisce i criteri, le metriche e gli standard da rispettare affinché il prodotto finale sia di alta qualità;
	\item \textbf{metodologie di testing}: contiene la descrizione di tutti i test necessari a verificare che il prodotto soddisfi i requisiti specificati;
	\item \href{https://7last.github.io/docs/pb/documentazione-interna/glossario\#cruscotto}{\textbf{cruscotto\textsubscript{G}} di valutazione della qualità}: riporta le attività di verifica svolte e le problematiche riscontrate durante lo sviluppo del software, con l'obiettivo di identificare aree di miglioramento;
	\item \textbf{iniziative di automiglioramento per la qualità}: contiene le azioni correttive e preventive adottate per migliorare la qualità del prodotto e dei processi.
\end{itemize}

\subsubsubsection{Analisi dei requisiti}
Questo documento, redatto dall'\href{https://7last.github.io/docs/pb/documentazione-interna/glossario\#analista}{analista\textsubscript{G}}, offre una panoramica chiara delle richieste e delle aspettative dell'azienda \href{https://7last.github.io/docs/pb/documentazione-interna/glossario\#proponente}{proponente\textsubscript{G}}. Serve come guida per il team di sviluppo, fornendo un elenco dettagliato delle funzionalità da implementare e dei requisiti da soddisfare. Infine, durante le fasi di verifica e validazione, questo documento viene utilizzato come base di riferimento per garantire che il prodotto software sviluppato rispetti le richieste e le necessità degli utenti finali, contribuendo così al successo complessivo del progetto.
Il documento è composto principalmente dalle sezioni qui specificate.
\begin{itemize}
	\item \textbf{Descrizione}: fornisce una panoramica generale del progetto, del contesto in cui il software verrà utilizzato e degli obiettivi principali. Include una spiegazione chiara dello scopo del documento e delle sue finalità.
	\item \textbf{Elenco dei casi d'uso}: riporta tutti gli scenari possibili in cui il sistema software potrebbe essere utilizzato dagli utenti finali. Ogni caso d'uso descrive dettagliatamente le azioni che gli utenti compiono nel sistema, permettendo di identificare requisiti non ovvi inizialmente. Questa sezione fornisce una visione dettagliata delle interazioni previste tra gli utenti e il sistema, contribuendo a definire le funzionalità e le caratteristiche del software.
	\item \textbf{Elenco dei requisiti}: specifica tutti i vincoli richiesti dalla \href{https://7last.github.io/docs/pb/documentazione-interna/glossario\#proponente}{proponente\textsubscript{G}} o dedotti in base all'analisi dei casi d'uso associati ad essi. I requisiti delineano le funzionalità, le prestazioni, le restrizioni e altri aspetti critici del software che devono essere soddisfatti. Essi rappresentano le basi su cui il team progetta e implementa il software, assicurando che il prodotto finale rispetti le aspettative degli \href{https://7last.github.io/docs/pb/documentazione-interna/glossario\#stakeholder}{stakeholder\textsubscript{G}} e soddisfi le necessità degli utenti finali.
\end{itemize}

\subsubsubsection{Piano di progetto}
Redatto dal \href{https://7last.github.io/docs/pb/documentazione-interna/glossario\#responsabile}{responsabile\textsubscript{G}}, delinea in modo dettagliato e organizzato tutti gli aspetti e le attività coinvolte nella gestione di un progetto di sviluppo del software. Esso funge da strumento di pianificazione, monitoraggio e controllo, offrendo una visione chiara e ben definita per il team di progetto. Il documento è costituito principalmente dalle sezioni qui elencate.
\begin{itemize}
	\item \textbf{Analisi dei rischi}: identificazione delle problematiche che potrebbero insorgere durante il processo e che potrebbero arrecare ritardi o impedimenti nella progressione del progetto. Il gruppo si impegna a fornire soluzioni per tali problemi il prima possibile. I rischi sono classificati in tre categorie principali: rischi organizzativi, rischi tecnologici e rischi comunicativi.
	\item \textbf{Calendario di progetto}: definisce le tempistiche pianificate per le varie revisioni previste durante lo svolgimento del progetto.
	\item \textbf{Stima costi realizzazione}: espone una previsione dettagliata dei costi di realizzazione del progetto.
	\item \textbf{Gestione del modello e \href{https://7last.github.io/docs/pb/documentazione-interna/glossario\#sprint}{sprint\textsubscript{G}}}: definisce i periodi temporali con eventi e attività correlate, all'interno di un calendario. Per ciascun periodo vengono analizzati i rischi, i costi e le attività sia in fase di pianificazione che in fase di retrospettiva.
	\item \textbf{Retrospettiva generale}: analisi generale di tutto il percorso svolto, ottenuta dall'esperienza accumulata, con particolare attenzione agli aspetti positivi e negativi riscontrati.
\end{itemize}


\subsubsubsection{Glossario}
Il \href{https://7last.github.io/docs/pb/documentazione-interna/glossario\#glossario}{\textit{Glossario}\textsubscript{G}} è un elenco dettagliato e organizzato di termini, acronimi e definizioni utilizzati nel contesto del progetto. Fornisce una chiara comprensione dei concetti e dei termini specifici impiegati nel progetto, garantendo una comunicazione efficace e coesa tra tutti i membri del team, nonché con gli \href{https://7last.github.io/docs/pb/documentazione-interna/glossario\#stakeholder}{stakeholder\textsubscript{G}}.

\subsubsubsection{Lettera di presentazione}
La \textit{Lettera di Presentazione} costituisce il mezzo attraverso il quale il gruppo \textit{7Last} espone il proprio interesse a partecipare attivamente alla fase di revisione del prodotto software. Questo documento non solo mette in luce la disponibilità della documentazione rilevante per i \href{https://7last.github.io/docs/pb/documentazione-interna/glossario\#committente}{committenti\textsubscript{G}} e la \href{https://7last.github.io/docs/pb/documentazione-interna/glossario\#proponente}{proponente\textsubscript{G}}, ma stabilisce anche i termini concordati per la consegna del prodotto finito.

\subsubsection{Metriche}
\begin{table}[!h]
	\centering
	\begin{tabular}{ | c | l | }
		\hline
		\textbf{Codice}                       & \textbf{Nome esteso}       \\
		\hline
		\underline{\hyperlink{1M}{1M-PV}}     & Planned Value              \\
		\underline{\hyperlink{2M}{2M-EV}}     & Earned Value               \\
		\underline{\hyperlink{3M}{3M-AC}}     & Actual Cost                \\
		\underline{\hyperlink{4M}{4M-SV}}     & Schedule Variance          \\
		\underline{\hyperlink{5M}{5M-CV}}     & Cost Variance              \\
		\underline{\hyperlink{6M}{6M-CPI}}    & Cost Performance Index     \\
		\underline{\hyperlink{7M}{7M-SPI}}    & Schedule Performance Index \\
		\underline{\hyperlink{8M}{8M-EAC}}    & Estimate at Completion     \\
		\underline{\hyperlink{9M}{9M-ETC}}    & Estimate to Complete       \\
		\underline{\hyperlink{10M}{10M-OTDR}} & On Time Delivery Rate      \\
		\hline
	\end{tabular}
	\caption{Metriche inerenti il processo di fornitura}
\end{table}

\subsubsection{Strumenti}
Di seguito sono descritti gli strumenti software impiegati nel processo di fornitura:
\begin{itemize}
	\item \textbf{Discord} come piattaforma per le riunioni interne e metodo informale per contattare l'azienda \href{https://7last.github.io/docs/pb/documentazione-interna/glossario\#proponente}{proponente\textsubscript{G}} tramite messaggistica;
	\item \textbf{Google Meet} utilizzato per le riunioni con la \href{https://7last.github.io/docs/pb/documentazione-interna/glossario\#proponente}{proponente\textsubscript{G}};
	\item \textbf{Google Slides} come strumento per la creazione di presentazioni;
	\item \textbf{LaTeX} come sistema di preparazione di documenti utilizzato principalmente per la creazione di documenti tecnici e scientifici;
	\item \href{https://7last.github.io/docs/pb/documentazione-interna/glossario\#clickup}{\textbf{ClickUp}\textsubscript{G}} come strumento di gestione del progetto per la pianificazione e il monitoraggio delle attività.
\end{itemize}

\subsection{Sviluppo}
\subsubsection{Introduzione}
Lo standard \textit{ISO/IEC 12207:1995} fornisce un quadro completo e strutturato per la gestione del ciclo di vita del software. Definisce i processi, le attività e le mansioni coinvolte nello sviluppo, nell'acquisizione e nella manutenzione del software.
È fondamentale completare tali operazioni in stretto rispetto alle direttive e ai requisiti definiti nel contratto con il cliente, assicurando quindi una realizzazione precisa e in linea con le specifiche richieste.

\subsubsection{Analisi dei requisiti}
\subsubsubsection{Descrizione}
Ha lo scopo di identificare, definire e documentare in modo esaustivo le necessità, le funzionalità e le prestazioni che il sistema software deve soddisfare. Questo documento rappresenta il punto di partenza per il processo di sviluppo del software, fornendo una base solida e chiara per la progettazione, l'implementazione e la verifica del sistema. Attraverso l'\href{https://7last.github.io/docs/pb/documentazione-interna/glossario\#analisi-dei-requisiti}{analisi dei requisiti\textsubscript{G}}, si cerca di comprendere appieno le esigenze degli \href{https://7last.github.io/docs/pb/documentazione-interna/glossario\#stakeholder}{stakeholder\textsubscript{G}}, inclusi utenti finali e clienti, al fine di garantire che il prodotto software sviluppato soddisfi le loro aspettative e necessità. L'\href{https://7last.github.io/docs/pb/documentazione-interna/glossario\#analisi-dei-requisiti}{analisi dei requisiti\textsubscript{G}} include tipicamente la raccolta e la documentazione dei requisiti funzionali di vincolo e qualitativi, la definizione dei casi d'uso nonché la prioritizzazione e la tracciabilità dei requisiti lungo l'intero ciclo di vita del software.
\newpage
\subsubsubsection{Scopo}
Questo processo pone i seguenti obiettivi principali:
\begin{itemize}
	\item promuovere una comprensione condivisa tra tutte le parti interessate;
	\item consentire una stima accurata delle tempistiche e dei costi del progetto;
	\item fornire ai \href{https://7last.github.io/docs/pb/documentazione-interna/glossario\#progettista}{progettisti\textsubscript{G}} requisiti chiari e facilmente comprensibili;
	\item definire gli obiettivi del prodotto al fine di soddisfare le aspettative;
	\item agevolare l'attività di verifica e di test fornendo indicazioni pratiche di riferimento.
\end{itemize}
\subsubsubsection{Documentazione}
Gli \href{https://7last.github.io/docs/pb/documentazione-interna/glossario\#analista}{analisti\textsubscript{G}} hanno il compito di redigere l'\href{https://7last.github.io/docs/pb/documentazione-interna/glossario\#analisi-dei-requisiti}{\textit{Analisi dei Requisiti}\textsubscript{G}}, comprendendo le seguenti sezioni:
\begin{itemize}
	\item \textbf{introduzione}, contente la presentazione e lo scopo del documento stesso;
	\item \textbf{descrizione del prodotto}, in cui è definita un'analisi approfondita del prodotto, includendo:
	      \begin{itemize}
		      \item obiettivi del prodotto;
		      \item architettura del prodotto;
		      \item funzionalità del prodotto;
		      \item caratteristiche degli utenti;
	      \end{itemize}
	\item \textbf{casi d'uso} con le descrizioni delle funzionalità offerte dal sistema dal punto di vista dell'utente, includendo:
	      \begin{itemize}
		      \item attori coinvolti;
		      \item elenco dei casi d'uso
	      \end{itemize}
	\item \textbf{requisiti}, suddivisi in:
	      \begin{itemize}
		      \item funzionali;
		      \item qualitativi;
		      \item di vincolo;
		      \item prestazionali.
	      \end{itemize}
\end{itemize}

\subsubsubsection{Casi d'uso}
I casi d'uso forniscono una descrizione dettagliata delle funzionalità del sistema dal punto di vista degli utenti, delineando come il sistema risponda a specifiche azioni o scenari. Ogni caso d'uso deve includere:
\begin{itemize}
	\item \textbf{identificativo}:
	    \begin{center}
		    \textbf{UC-[identificativo\_caso\_principale].[identificativo\_sotto\_caso]}
	    \end{center}
	    \begin{itemize}
		    \item \textbf{identificativo caso principale}: identificativo numerico del caso d'uso principale;
		    \item \textbf{identificativo sotto caso}: identificativo numerico del sotto caso d'uso (presente solo se si tratta di un sotto caso d'uso);
		\end{itemize}
	\item \textbf{titolo}: breve e chiaro titolo del caso d'uso;
	\item \textbf{attore principale}: entità esterna che interagisce attivamente con il sistema per soddisfare una propria necessità;
	\item \textbf{precondizioni}: lo stato in cui deve trovarsi il sistema affinché la funzionalità sia disponibile per l'attore;
	\item \textbf{postcondizioni}: lo stato in cui si trova il sistema dopo l'esecuzione dello scenario principale;
	\item \textbf{scenario principale}: una sequenza di eventi che si verificano quando un attore interagisce con il sistema per raggiungere l'obiettivo del caso d'uso (postcondizioni);
	\item \href{https://7last.github.io/docs/pb/documentazione-interna/glossario\#user-story}{\textbf{user story}\textsubscript{G}}: una breve descrizione di una funzionalità del software, scritta dal punto di vista dell'utente, che fornisce contesto, obiettivi e valore:
	      \begin{itemize}
		      \item L'\href{https://7last.github.io/docs/pb/documentazione-interna/glossario\#user-story}{user story\textsubscript{G}} viene scritta nella forma: "Come [utente] desidero poter [funzionalità] per [valore aggiunto]";
	      \end{itemize}
\end{itemize}

\subsubsubsection{Diagrammi dei casi d'uso}
Offrono una rappresentazione visiva delle interazioni tra gli attori e il sistema, delineando i vari scenari in cui esso verrà utilizzato. Ogni caso d'uso descrive una sequenza specifica di azioni che gli attori compiono per raggiungere un obiettivo particolare all'interno del sistema. I diagrammi dei casi d'uso aiutano a identificare i requisiti funzionali, fornendo una comprensione chiara e condivisa delle aspettative degli utenti. Inoltre, facilitano la comunicazione tra gli sviluppatori e gli \href{https://7last.github.io/docs/pb/documentazione-interna/glossario\#stakeholder}{stakeholder\textsubscript{G}}, assicurando che tutte le funzionalità necessarie siano considerate e correttamente implementate nell'applicativo.
Di seguito vengono elencati i principali componenti di un diagramma dei casi d'uso.
\begin{itemize}
	\item \textbf{Sistema}: delimita i confini del sistema software, indicando quali funzionalità sono incluse e quali sono esterne ad esso.
	      \begin{center}
		      \includegraphics*[width=8cm]{../../../images/norme_di_progetto/sistema.png}
		      \captionof{figure}{Diagramma dei casi d'uso - sistema}
	      \end{center}
	\item \textbf{Attori}: rappresentano soggetti che interagiscono con il sistema software. Gli attori possono essere persone, altri applicativi o dispositivi che utilizzano le funzionalità del sistema.
	      \begin{center}
		      \includegraphics*[width=2.5cm]{../../../images/norme_di_progetto/attore.png}
		      \captionof{figure}{Diagramma dei casi d'uso - attore}
	      \end{center}
	\item \textbf{Casi d'uso}: rappresentano le funzionalità o i servizi offerti dal sistema che producono un risultato di valore per un attore. Ogni caso d'uso descrive una sequenza specifica di interazioni tra gli attori e il sistema.
	      \begin{center}
		      \includegraphics*[width=6cm]{../../../images/norme_di_progetto/uc.png}
		      \captionof{figure}{Diagramma dei casi d'uso - caso d'uso}
	      \end{center}
	\item \textbf{Sottocasi d'uso}: rappresentano scenari specifici e dettagliati che si verificano all'interno di un caso d'uso principale, descrivendolo in modo più dettagliato.
	      \begin{center}
		      \includegraphics*[width=8cm]{../../../images/norme_di_progetto/sottocasi.png}
		      \captionof{figure}{Diagramma dei casi d'uso - sottocaso d'uso}
	      \end{center}

	\item \textbf{Relazioni tra Attori e Casi d'Uso}
	      \begin{itemize}
		      \item \textbf{Associazione}: è la relazione fondamentale che collega un attore a un caso d'uso, indica quali attori interagiscono con quali casi d'uso.
		            \begin{center}
			            \includegraphics*[width=8cm]{../../../images/norme_di_progetto/associazione.png}
			            \captionof{figure}{Diagramma dei casi d'uso - associazione}
		            \end{center}
	      \end{itemize}

	\item \textbf{Relazioni tra Attori}
	      \begin{itemize}
		      \item \textbf{Generalizzazione}: rappresenta una relazione in cui un attore figlio condivide almeno le funzionalità di un attore genitore. Utilizzata quando diversi attori condividono comportamenti comuni.
		            \begin{center}
			            \includegraphics*[width=1.5cm]{../../../images/norme_di_progetto/generalizzazione_attori.png}
			            \captionof{figure}{Diagramma dei casi d'uso - generalizzazione tra attori}
		            \end{center}
	      \end{itemize}

	\item \textbf{Relazioni tra Casi d'Uso}
	      \begin{itemize}
		      \item \textbf{Inclusione}: rappresenta una dipendenza in cui il comportamento del caso d'uso incluso è incorporato ogni volta che viene eseguito il caso d'uso base. Il caso d'uso incluso contiene funzionalità che sono riutilizzate in più casi d'uso principali, permettendo di evitare la duplicazione di comportamenti comuni.
		            \begin{center}
			            \includegraphics*[width=7cm]{../../../images/norme_di_progetto/include.png}
			            \captionof{figure}{Diagramma dei casi d'uso - inclusione}
		            \end{center}
		      \item \textbf{Estensione}: mostra una relazione in cui un caso d'uso esteso aumenta le funzionalità del caso d'uso principale. Il caso d'uso esteso viene eseguito solo sotto determinate condizioni, interrompendo l'esecuzione del caso d'uso principale.
		            \begin{center}
			            \includegraphics*[width=8cm]{../../../images/norme_di_progetto/estensione.png}
			            \captionof{figure}{Diagramma dei casi d'uso - estensione}
		            \end{center}
		      \item \textbf{Generalizzazione}: rappresenta una relazione in cui un caso d'uso figlio può aggiungere funzionalità o modificare il comportamento di un caso d'uso genitore. Tutte le funzionalità definite nel caso d'uso genitore si mantengono nel caso d'uso figlio se queste non vengono ridefinite.
		            \begin{center}
			            \includegraphics*[width=9cm]{../../../images/norme_di_progetto/generalizzazione_uc.png}
			            \captionof{figure}{Diagramma dei casi d'uso - generalizzazione tra casi d'uso}
		            \end{center}
	      \end{itemize}


\end{itemize}

\subsubsubsection{Requisiti}
Delineano una descrizione dettagliata delle funzionalità e delle caratteristiche che il software deve possedere per soddisfare le esigenze degli utenti e degli \href{https://7last.github.io/docs/pb/documentazione-interna/glossario\#stakeholder}{stakeholder\textsubscript{G}}. Questi requisiti costituiscono la base per tutto il ciclo di sviluppo del software e sono fondamentali per il successo del progetto. Si suddividono generalmente in tre categorie principali:
\begin{itemize}
	\item \textbf{requisiti funzionali}, i quali descrivono le azioni specifiche che il sistema deve essere in grado di eseguire;
	\item \textbf{requisiti non funzionali}, che definiscono le proprietà e le caratteristiche del sistema che ne influenzano il funzionamento e l'uso;
	\item \textbf{requisiti di vincolo}, che rappresentano i vincoli e le restrizioni che il sistema deve rispettare.
\end{itemize}
Ogni requisito è costituito da:
\begin{enumerate}
	\item \textbf{codice} - i requisiti sono codificati nel seguente modo:
	      \begin{center}
		      \textbf{R[tipologia]-[identificativo]}
	      \end{center}
	      dove:
	      \begin{itemize}
		      \item \textbf{[tipologia]} indica la tipologia del requisito, che può essere
		            \begin{itemize}
			            \item \textbf{F} = requisito funzionale
			            \item \textbf{Q} = requisito di qualità
			            \item \textbf{V} = requisito di vincolo;
		            \end{itemize}
		      \item \textbf{[identificativo]} è un numero progressivo che identifica univocamente il requisito;
	      \end{itemize}
	\item \textbf{importanza} - indica il grado di importanza del requisito, che può essere:
	      \begin{itemize}
		      \item \textbf{obbligatorio}, ovvero irrinunciabile per il \href{https://7last.github.io/docs/pb/documentazione-interna/glossario\#committente}{committente\textsubscript{G}};
		      \item \textbf{desiderabile} cioè non strettamente necessario, ma che porta valore aggiunto al prodotto;
		      \item \textbf{opzionale}, relativo a funzionalità aggiuntive;
	      \end{itemize}
	\item \textbf{fonte} - indica la fonte da cui è stato identificato il requisito, che può essere:
	      \begin{itemize}
		      \item \href{https://7last.github.io/docs/pb/documentazione-interna/glossario\#capitolato}{\textbf{capitolato}\textsubscript{G}} indica i requisiti individuati a seguito dell'analisi dello stesso;
		      \item \textbf{interno} identifica i requisiti individuati durante le riunioni interne e da coloro che hanno il ruolo di \href{https://7last.github.io/docs/pb/documentazione-interna/glossario\#analista}{analista\textsubscript{G}};
		      \item \textbf{esterno} segnala i requisiti aggiuntivi individuati in seguito a incontri con la \href{https://7last.github.io/docs/pb/documentazione-interna/glossario\#proponente}{proponente\textsubscript{G}};
		      \item \href{https://7last.github.io/docs/pb/documentazione-interna/glossario\#piano-di-qualifica}{\textbf{piano di qualifica}\textsubscript{G}} evidenzia i requisiti necessari per adeguare il prodotto agli standard di qualità definiti nel documento \href{https://7last.github.io/docs/pb/documentazione-interna/glossario\#piano-di-qualifica}{\textit{Piano di Qualifica}\textsubscript{G}};
		      \item \href{https://7last.github.io/docs/pb/documentazione-interna/glossario\#norme-di-progetto}{\textbf{norme di progetto}\textsubscript{G}} specifica i requisiti necessari per adeguare il prodotto alle norme stabilite nel documento \href{https://7last.github.io/docs/pb/documentazione-interna/glossario\#norme-di-progetto}{\textit{Norme di Progetto}\textsubscript{G}}.
	      \end{itemize}
\end{enumerate}

\subsubsubsection{Metriche}
Nell'\href{https://7last.github.io/docs/pb/documentazione-interna/glossario\#analisi-dei-requisiti}{analisi dei requisiti\textsubscript{G}}, le metriche sono strumenti utilizzati per misurare vari aspetti del processo di sviluppo e del prodotto finale. All'interno di un'\href{https://7last.github.io/docs/pb/documentazione-interna/glossario\#analisi-dei-requisiti}{analisi dei requisiti\textsubscript{G}}, le metriche svolgono un ruolo cruciale nel garantire che il software sia sviluppato in modo efficiente, soddisfi gli standard di qualità attesi e adempia alle specifiche richieste dagli \href{https://7last.github.io/docs/pb/documentazione-interna/glossario\#stakeholder}{stakeholder\textsubscript{G}}.
\begin{table}[!h]
	\centering
	\begin{tabular}{ | c | l | }
		\hline
		\textbf{Codice}                      & \textbf{Nome esteso}               \\
		\hline
		\underline{\hyperlink{11M}{11M-PRO}} & Percentuale requisiti obbligatori  \\
		\underline{\hyperlink{12M}{12M-PRD}} & Percentuale requisiti desiderabili \\
		\underline{\hyperlink{13M}{13M-PRO}} & Percentuale requisiti opzionali    \\
		\hline
	\end{tabular}
	\caption{Metriche inerenti i requisiti}
\end{table}

\newpage

\subsubsubsection{Strumenti}
\textbf{Draw.io} è uno strumento online che permette di creare diagrammi dei casi d'uso in modo semplice e intuitivo. Questo strumento offre una vasta gamma di funzionalità e opzioni di personalizzazione per creare diagrammi chiari e ben strutturati. Draw.io è un'applicazione web-based, quindi non richiede alcun download o installazione, rendendo facile e veloce la creazione di diagrammi dei casi d'uso.

\subsubsection{Progettazione}
\subsubsubsection{Descrizione}
L'attività di progettazione svolge un ruolo fondamentale nello sviluppo del software, consentendo di definire l'architettura, i componenti e le interfacce del sistema in modo da soddisfare i requisiti definiti durante l'analisi.  Questa fase si basa sull'\href{https://7last.github.io/docs/pb/documentazione-interna/glossario\#analisi-dei-requisiti}{analisi dei requisiti\textsubscript{G}} e ha lo scopo di tradurre i requisiti funzionali e non funzionali in una soluzione tecnica implementabile.

\subsubsubsection{Obiettivi}
L'obiettivo principale è soddisfare i requisiti attraverso un sistema di qualità definito dall'architettura del prodotto.
\begin{itemize}
	\item \textbf{Definizione dell'architettura}: in questa fase si definisce l'architettura del sistema, determinando la struttura generale, i componenti principali e le relazioni tra di essi. Questo include la scelta delle tecnologie, dei linguaggi di programmazione e delle piattaforme da utilizzare.
	\item \textbf{Suddivisione in moduli}: si suddivide il sistema in moduli o componenti più piccoli e gestibili, ciascuno con responsabilità specifiche. Questa suddivisione facilita lo sviluppo, il testing e la manutenzione del software.
	\item \textbf{Definizione delle interfacce}: si specificano le interfacce tra i diversi moduli o componenti del sistema, stabilendo come essi comunicano e interagiscono tra di loro. Questo garantisce una corretta integrazione e interoperabilità del sistema.
	\item \textbf{Identificazione delle risorse}: vengono identificate le risorse necessarie per implementare il sistema, come hardware, software, risorse umane e finanziarie. Questo aiuta a pianificare e gestire le risorse in modo efficace durante lo sviluppo.
	\item \textbf{Valutazione delle prestazioni}: si valutano le prestazioni del sistema, identificando potenziali problemi di scalabilità, affidabilità e prestazioni. Questo consente di ottimizzare il design per garantire che il sistema soddisfi i requisiti di prestazione.
\end{itemize}

\begin{flushleft}
	Lo sviluppo del progetto inizia con l'analisi delle tecnologie suggerite dalla \href{https://7last.github.io/docs/pb/documentazione-interna/glossario\#proponente}{proponente\textsubscript{G}} e un confronto con eventuali alternative, valutando i pro e i contro di ciascuna. Dopo aver selezionato le tecnologie più adatte, si procede con lo sviluppo di un \href{https://7last.github.io/docs/pb/documentazione-interna/glossario\#proof-of-concept}{Proof of Concept\textsubscript{G}} (\href{https://7last.github.io/docs/pb/documentazione-interna/glossario\#proof-of-concept}{PoC\textsubscript{G}}), un prototipo grezzo che dimostra l'idoneità delle tecnologie scelte e la capacità del team di utilizzarle efficacemente. Questo \href{https://7last.github.io/docs/pb/documentazione-interna/glossario\#proof-of-concept}{PoC\textsubscript{G}} costituisce la componente principale della Technology Baseline, che funge da base per lo sviluppo del prodotto finale. A partire da questa, si procede con lo sviluppo del prodotto vero e proprio, trasformando il \href{https://7last.github.io/docs/pb/documentazione-interna/glossario\#proof-of-concept}{PoC\textsubscript{G}} in un \href{https://7last.github.io/docs/pb/documentazione-interna/glossario\#minimum-viable-product}{Minimum Viable Product\textsubscript{G}} (\href{https://7last.github.io/docs/pb/documentazione-interna/glossario\#minimum-viable-product}{MVP\textsubscript{G}}), ovvero un prodotto minimale che soddisfa i requisiti minimi richiesti e che diventerà parte integrante della \href{https://7last.github.io/docs/pb/documentazione-interna/glossario\#product-baseline}{Product Baseline\textsubscript{G}} (\href{https://7last.github.io/docs/pb/documentazione-interna/glossario\#product-baseline}{PB\textsubscript{G}}).
\end{flushleft}

\subsubsubsection{Documentazione}
\textbf{Specifica tecnica}\\
Fornisce una guida dettagliata ed esaustiva per la progettazione e lo sviluppo del sistema. In essa sono delineate con precisione l'architettura tecnica del software, con tutti i suoi componenti e interfacce, le funzionalità che il sistema dovrà offrire e le tecnologie che saranno impiegate per realizzarlo. Questo documento costituisce il punto di partenza per l'intero ciclo di sviluppo del software, fornendo una base solida e strutturata su cui il team di sviluppo potrà lavorare in modo efficace ed efficiente. Tra gli elementi chiave inclusi in questo documento vi sono:
\begin{itemize}
	\item \textbf{architettura del sistema}: questa sezione delinea l'architettura generale del sistema software, identificando i principali componenti, i moduli e le interfacce. Include anche informazioni sull'organizzazione logica e fisica del sistema, come ad esempio la suddivisione in livelli o strati e le relazioni tra di essi;
	\item \textbf{tecnologie utilizzate}: specifica le tecnologie, i linguaggi di programmazione, i framework e le librerie che saranno impiegate nello sviluppo del sistema. Questo include anche eventuali tool e ambienti di sviluppo adottati;
	\item \textbf{struttura dei dati}: descrive la struttura e la gestione dei dati all'interno del sistema, inclusi database, file system, protocolli di accesso ai dati e modelli di persistenza;
	\item \textbf{interfacce esterne}: indica le interfacce con altri sistemi o applicazioni esterne con cui il sistema deve interagire. Questo include anche la definizione dei formati dei dati scambiati e i protocolli di comunicazione utilizzati;
	\item \textbf{design pattern}: descrizione dei design pattern utilizzati per risolvere problemi comuni e ricorrenti durante lo sviluppo del software;
	\item \textbf{pianificazione e risorse}: fornisce una pianificazione dettagliata del lavoro da svolgere per implementare la specifica tecnica, inclusi tempi, costi e risorse necessarie;
	\item \textbf{procedure di testing e validazione}: indicazioni sulle procedure e gli strumenti da utilizzare per verificare e validare il prodotto, garantendo che soddisfi i requisiti e le aspettative del cliente;
	\item \textbf{requisiti tecnici}: elenco dettagliato dei requisiti tecnici che il prodotto deve soddisfare, con indicazioni sulle funzionalità, le prestazioni e le caratteristiche richieste;
\end{itemize}

\subsubsubsection{Qualità dell'architettura}
\begin{itemize}
	\item \textbf{Scalabilità}: rappresenta la capacità del sistema di gestire un aumento del carico di lavoro o delle risorse senza compromettere le prestazioni o la qualità del servizio. Un'architettura scalabile può adattarsi dinamicamente alle variazioni delle richieste degli utenti e delle risorse disponibili.
	\item \textbf{Flessibilità}: capacità del sistema di adattarsi ai cambiamenti nei requisiti o nell'ambiente operativo senza richiedere modifiche significative. Un'architettura flessibile è caratterizzata da componenti ben separati e interfacce standardizzate, che consentono una facile modifica e l'aggiunta di nuove funzionalità.
	\item \textbf{Manutenibilità}: ovvero la facilità di comprendere, modificare e correggere il software nel corso del suo ciclo di vita. Un'architettura manutenibile è caratterizzata da un codice ben strutturato, documentato e modularizzato, che consente agli sviluppatori di individuare e risolvere rapidamente eventuali problemi.
	\item \textbf{Testabilità}: facilità di testare il software per garantire che soddisfi i requisiti funzionali e non funzionali. Un'architettura testabile include componenti ben isolati e interfacce chiare che consentono l'automazione dei test e la verifica continua della qualità del software.
	\item \textbf{Affidabilità}: garanzia del corretto funzionamento del sistema in condizioni normali e anomale. Un'architettura affidabile include meccanismi di gestione degli errori, di recupero e di ripristino che consentono al sistema di continuare a operare in modo affidabile anche in presenza di guasti o interruzioni.
	\item \textbf{Performance}: capacità del sistema di fornire risposte rapide e tempi di elaborazione efficienti anche in presenza di carichi di lavoro elevati. Un'architettura performante include ottimizzazioni del codice, della gestione delle risorse e dell'accesso ai dati per massimizzare le prestazioni del sistema.
	\item \textbf{Usabilità}: facilità di utilizzo e comprensione del sistema da parte degli utenti finali. Un'architettura usabile include un'interfaccia utente intuitiva, una struttura di navigazione chiara e una presentazione coerente delle informazioni per garantire un'esperienza utente positiva.
	\item \textbf{Compatibilità}: capacità del sistema di interoperare con altri sistemi, piattaforme o tecnologie senza problemi. Un'architettura compatibile include standard aperti, interfacce ben definite e protocolli di comunicazione standardizzati che consentono l'integrazione con sistemi esterni.
	\item \textbf{Documentazione}: fornitura di documentazione completa e accurata sull'architettura del sistema, inclusi diagrammi, specifiche tecniche, manuale utente e guide per gli sviluppatori. Una buona documentazione facilita la comprensione, la manutenzione e l'evoluzione del sistema nel tempo.
	\item \textbf{Safety}: l'architettura deve garantire la sicurezza del sistema, in modo che possa proteggere i dati e le informazioni sensibili in seguito a malfunzionamenti.
\end{itemize}

\subsubsubsection{Diagrammi UML}
\begin{itemize}
	\item \textbf{Chiarezza visiva}: forniscono una rappresentazione visuale chiara e intuitiva delle relazioni e delle interazioni tra gli elementi del sistema software, facilitando la comprensione del sistema da parte degli \href{https://7last.github.io/docs/pb/documentazione-interna/glossario\#stakeholder}{stakeholder\textsubscript{G}}.
	\item \textbf{Standardizzazione}: UML è uno standard riconosciuto a livello internazionale per la modellazione dei sistemi software, il che significa che i diagrammi UML sono facilmente comprensibili da parte di professionisti del settore in tutto il mondo.
	\item \textbf{Comunicazione efficace}: forniscono un linguaggio comune per la comunicazione tra gli sviluppatori, gli \href{https://7last.github.io/docs/pb/documentazione-interna/glossario\#stakeholder}{stakeholder\textsubscript{G}} tecnici e non tecnici e altri membri del team, consentendo una comunicazione più efficace e una condivisione delle informazioni più chiara.
	\item \textbf{Analisi e progettazione}: possono essere utilizzati durante le fasi di analisi e progettazione del ciclo di sviluppo del software per visualizzare e analizzare i requisiti, le relazioni tra i componenti e le interazioni del sistema.
	\item \textbf{Modellazione e requisiti}: UML consente di modellare i requisiti del sistema in modo chiaro e strutturato, identificando casi d'uso, scenari e vincoli che guidano lo sviluppo del software.
	\item \textbf{Facilità di manutenzione}: è più facile comprendere l'architettura del sistema e individuare eventuali problemi o aree di miglioramento, facilitando così la manutenzione e l'evoluzione del software nel tempo.
	\item \textbf{Supporto agli standard di progettazione}: i diagrammi UML possono essere utilizzati per applicare e comunicare i principi di progettazione software ben consolidati, come l'incapsulamento, l'ereditarietà e il polimorfismo.
	\item \textbf{Documentazione}: possono essere utilizzati per generare documentazione tecnica dettagliata sul sistema software, fornendo una guida completa per gli sviluppatori, gli utenti finali e gli altri \href{https://7last.github.io/docs/pb/documentazione-interna/glossario\#stakeholder}{stakeholder\textsubscript{G}}.
\end{itemize}

\subsubsubsubsection{Diagrammi delle Classi}
Un tipo di diagramma UML ampiamente utilizzato nel processo di progettazione del software è il diagramma delle classi, che rappresenta la struttura statica del sistema, mostrando le classi, gli attributi, i metodi e le relazioni tra di esse. Le classi sono rappresentate da rettangoli suddivisi in tre sezioni:
\begin{enumerate}
	\item \textbf{nome della classe} - indica il nome della classe;
	\item \textbf{attributi} - elenco degli attributi della classe, con il relativo tipo di dato, seguendo il formato: 
		\begin{center}
			\textbf{visibilità nome: tipo [molteplicità] = default}
		\end{center}
		dove:
		\begin{itemize}
			\item \textbf{visibilità} indica il livello di accesso agli attributi, che può essere
				\begin{itemize}
					\item \textbf{+} = pubblico
					\item \textbf{-} = privato
					\item \textbf{\#} = protetto
					\item \textbf{\textasciitilde} = package;
				\end{itemize}
			\item \textbf{nome} indica il nome dell'attributo che deve essere rappresentativo, chiaro e deve seguire la notazione \textit{nomeAttributo: tipo} e se l'attributo è costante il nome deve essere scritto in maiuscolo (es. \textit{PIGRECO: double});
			\item \textbf{Molteplicità} presente nel caso di una sequenza di elementi come liste o array, indica il numero di elementi presenti, se questo non fosse conosciuto si utilizza il simbolo \textit{*} (es \textit{tipoAttributo[*]});
			\item \textbf{default} indica il valore di default dell'attributo;
		\end{itemize}
	\item \textbf{metodi} - descrivono il comportamento della classe, seguendo il formato:
		\begin{center}
			\textbf{visibilità nome(parametri): ritorno}
		\end{center}
		dove:
		\begin{itemize}
			\item \textbf{visibilità} indica il livello di accesso ai metodi, che può essere
				\begin{itemize}
					\item \textbf{+} = pubblico
					\item \textbf{-} = privato
					\item \textbf{\#} = protetto
					\item \textbf{\textasciitilde} = package;
				\end{itemize}
			\item \textbf{nome}, indica il nome del metodo, deve essere rappresentativo, chiaro e deve seguire la notazione \textit{nomeMetodo(parametri): tipoRitorno};
			\item \textbf{parametri} rappresenta l'elenco dei parametri del metodo, separati tramite virgola, dove ogni parametro deve seguire la notazione \textit{nomeParametro: tipo};
			\item \textbf{ritorno} indica il tipo di dato restituito dal metodo.
		\end{itemize}
\end{enumerate}

\newpage
\subsubsubsubsubsection*{Convenzioni sui metodi}
Le indicazioni dei metodi all'interno del diagramma delle classi seguono le seguenti convenzioni:
\begin{itemize}
	\item \textbf{i metodi getter, setter} e i \textbf{costruttori} non vengono inclusi fra i metodi;
	\item \textbf{i metodi statici} sono sottolineati;
	\item \textbf{i metodi astratti} sono scritti in corsivo;
	\item \textbf{l'assenza di attributi o metodi} in una classe determina l'assenza delle relative sezioni nel diagramma.
\end{itemize}
\subsubsubsubsubsection*{Relazioni tra le classi}
Essenziali per rappresentare le interazioni e le dipendenze tra le varie componenti del sistema software, queste relazioni forniscono un quadro completo della struttura e del comportamento del sistema, consentendo agli sviluppatori di comprendere meglio come le classi si relazionano tra loro e di progettare un sistema coerente e ben strutturato. Di seguito sono descritte le principali relazioni tra le classi:
\begin{itemize}
	\item \textbf{Associazione}:  rappresenta una relazione tra due classi, indicando che un'istanza di una classe è correlata a un'istanza di un'altra classe. È raffigurata da una linea solida tra le classi coinvolte. La linea può includere etichette per indicare il nome dell'associazione, nonché molteplicità e ruoli per specificare la cardinalità e il ruolo delle classi nell'associazione.
		\begin{center}
			\includegraphics*[width=12cm]{../../../images/norme_di_progetto/associazione_classi.png}
			\captionof{figure}{Diagramma delle classi - associazione}
		\end{center}
	\newpage
	\item \textbf{Aggregazione}:  indica una relazione "tutto-parte" tra due classi, dove una classe è composta da una o più istanze di un'altra classe, ma le istanze possono esistere indipendentemente dalla classe principale. Composta da una linea con una freccia vuota che punta dalla parte composta (parte) al tutto (oggetto principale). La linea può includere un rombo vuoto sulla parte composta per indicare l'aggregazione.
		\begin{center}
			\includegraphics*[width=12cm]{../../../images/norme_di_progetto/aggregazione_classi.png}
			\captionof{figure}{Diagramma delle classi - aggregazione}
		\end{center}
	\item \textbf{Composizione}: simile all'aggregazione, ma indica una relazione più forte in cui l'esistenza delle parti dipende direttamente dall'oggetto principale. Una composizione è rappresentata da una linea con una freccia piena che punta dalla parte composta (parte) al tutto (oggetto principale). La linea può includere un rombo pieno sulla parte composta per indicare la composizione.
		\begin{center}
			\includegraphics*[width=10cm]{../../../images/norme_di_progetto/composizione_classi.png}
			\captionof{figure}{Diagramma delle classi - composizione}
		\end{center}
	\newpage
	\item \textbf{Ereditarietà o generalizzazione}: indica una relazione di specializzazione tra una classe superiore (superclasse) e una o più classi inferiori (sottoclassi). Le sottoclassi ereditano gli attributi e i metodi della superclasse e possono estenderli o modificarli. Viene effigiata da una linea con una freccia che punta dalla sottoclasse alla superclasse. La linea può includere una freccia vuota con la scritta "extends" per indicare la generalizzazione.
		\begin{center}
			\includegraphics*[width=10cm]{../../../images/norme_di_progetto/generalizzazione_classi.png}
			\captionof{figure}{Diagramma delle classi - generalizzazione}
		\end{center}
	\item \textbf{Dipendenza}: indica che una classe dipende dall'altra in un certo modo. Questo può significare che una classe utilizza o richiede i servizi di un'altra classe, ma non c'è un legame diretto tra le loro istanze. Una dipendenza è rappresentata da una linea tratteggiata che parte dalla classe dipendente e punta alla classe di cui dipende. La linea può includere una freccia vuota per indicare la direzione della dipendenza.
		\begin{center}
			\includegraphics*[width=12cm]{../../../images/norme_di_progetto/dipendenza_classi.png}
			\captionof{figure}{Diagramma delle classi - dipendenza}
		\end{center}
	\item \textbf{Realizzazione}:  indica che una classe implementa un'interfaccia o un contratto definito da un'altra classe. In altre parole, una classe realizza i metodi definiti da un'interfaccia o una classe astratta. Una realizzazione è raffigurata da una linea tratteggiata con una freccia vuota che punta dalla classe che implementa all'interfaccia o classe astratta che viene implementata.
	      \begin{center}
		      \includegraphics*[width=8cm]{../../../images/norme_di_progetto/realizzazione_classi.png}
		      \captionof{figure}{Diagramma delle classi - realizzazione}
	      \end{center}
\end{itemize}

\subsubsubsection{Design pattern}
I design pattern  sono soluzioni progettuali generiche e riutilizzabili per problemi comuni che si presentano durante lo sviluppo del software. Sono modelli architetturali, concetti o paradigmi consolidati, spesso espressi attraverso codice o diagrammi, che forniscono una guida per risolvere specifiche problematiche di progettazione software in modo efficiente ed efficace. la documentazione sui design pattern serve come punto di riferimento centrale per tutti gli sviluppatori coinvolti nel progetto, consentendo loro di comprendere rapidamente e facilmente come sono stati implementati determinati pattern nel codice. Questo è particolarmente utile durante le fasi di progettazione, sviluppo e manutenzione del software, in quanto fornisce una chiara visione della struttura architetturale del sistema e dei pattern utilizzati.

\subsubsubsection{Test}
La parte dei test all'interno del processo di sviluppo software è un'attività critica volta a garantire che il prodotto finale soddisfi i requisiti funzionali, prestazionali e di qualità previsti. Questa fase coinvolge diverse attività, che includono la pianificazione dei test, la progettazione dei casi di test, l'esecuzione dei test e l'analisi dei risultati. In questa sezione \hyperref[testing]{\textit{3.2.4}} vengono fornite descrizioni dettagliate delle varie tipologie di test e della terminologia associata.

\subsubsubsection{Metriche}
\begin{table}[!h]
	\centering
	\begin{tabular}{ | c | l | }
		\hline
		\textbf{Codice}                     & \textbf{Nome esteso}       \\
		\hline
		\underline{\hyperlink{14M}{14M-PG}} & Profondità delle Gerarchie \\
		\hline
	\end{tabular}
	\caption{Metriche inerenti la progettazione}
\end{table}

\subsubsubsection{Strumenti}
\textbf{Draw.io}: è uno strumento per la creazione di diagrammi online gratuito che consente di creare, modificare e condividere diagrammi di flusso, mappe mentali, organigrammi e altri tipi di diagrammi grafici.

\subsubsection{Codifica}
\subsubsubsection{Descrizione}
La parte di codifica, nota anche come implementazione, rappresenta una fase critica e fondamentale nel processo di sviluppo del software. Durante questa fase, gli sviluppatori assumono il compito di tradurre i requisiti funzionali e non funzionali, delineati durante l' analisi e la progettazione, in codice eseguibile. Questo processo richiede non solo competenze tecniche avanzate, ma anche una solida comprensione degli obiettivi del progetto e delle esigenze degli \href{https://7last.github.io/docs/pb/documentazione-interna/glossario\#stakeholder}{stakeholder\textsubscript{G}}.

\subsubsubsection{Obiettivi}
L'obiettivo principale della codifica è trasformare in modo accurato e efficiente i concetti astratti dei requisiti in istruzioni precise e comprensibili per il computer. Ciò implica la scrittura di codice pulito, ben strutturato e facilmente manutenibile, che sia in grado di soddisfare le esigenze funzionali del software e rispettare gli standard di qualità e le linee guida del progetto.

\subsubsubsection{Norme di codifica}
Le seguenti norme sono state formalizzate in questa maniera:
\begin{itemize}
	\item \textbf{nomi significativi}: i nomi delle variabili, delle costanti, delle classi e dei metodi devono essere significativi e rappresentativi della loro funzione, in modo da facilitare la comprensione del codice;
	\item \textbf{commenti}: il codice deve essere implementato nel modo più leggibile ed autoesplicativo possibile. I commenti verranno impiegati solo ove strettamente necessari;
	\item \textbf{indentazione e formattazione consistente}: il codice deve essere correttamente indentato e formattato, mediante l'uso di tabulazioni, per garantire una corretta leggibilità e comprensione ;
	\item \textbf{gestione delle eccezioni}: il codice deve gestire correttamente le eccezioni e gli errori, fornendo messaggi significativi e gestendo le situazioni anomale in modo appropriato;
	\item \textbf{riuso del codice}: il codice deve essere scritto in modo modulare e riutilizzabile;
	\item \textbf{testabilità}: scrivere funzioni piccole e focalizzate, ad esempio, aiuta a garantire che il codice sia facilmente testabile;
	\item \textbf{sicurezza}: il codice deve essere scritto in modo sicuro, evitando vulnerabilità e falle di sicurezza;
	\item \textbf{performance}: il codice deve essere ottimizzato per garantire prestazioni elevate e tempi di risposta rapidi, con l'uso corretto delle risorse e l'ottimizzazione dei cicli di elaborazione;
	\item \textbf{compatibilità}: il codice deve essere compatibile con le diverse piattaforme, sistemi operativi e browser, per garantire un'esperienza utente uniforme e coerente;
	\item \textbf{conformità ai principi SOLID}: il codice deve rispettare i principi SOLID, che promuovono la scrittura di codice pulito, modulare e manutenibile.
\end{itemize}

\subsubsubsection{Metriche}
\begin{table}[!h]
	\centering
	\begin{tabular}{|c|l|}
		\hline
		\textbf{Codice}                       & \textbf{Nome esteso}       \\
		\hline
		\underline{\hyperlink{15M}{15M-PPM}}  & Parametri per Metodo       \\
		\underline{\hyperlink{16M}{16M-CPC}}  & Campi per classe           \\
		\underline{\hyperlink{17M}{17M-LCPM}} & Linee di codice per metodo \\
		\underline{\hyperlink{18M}{18M-CCM}}  & Complessità Ciclomatica    \\
		\hline
	\end{tabular}
	\caption{Metriche riguardanti l'attività di codifica}
\end{table}

\subsubsubsection{Strumenti}
\textbf{Visual Studio Code}: è un ambiente di sviluppo integrato (IDE) sviluppato da Microsoft e adottato dal gruppo per lo sviluppo del software. Offre funzionalità avanzate per la scrittura, la modifica e il debug del codice, con supporto per numerosi linguaggi di programmazione e framework.

\newpage
\subsubsection{Configurazione dell'ambiente di esecuzione}
\subsubsubsection{Docker}
L'utilizzo del file \href{https://7last.github.io/docs/pb/documentazione-interna/glossario\#docker}{Docker\textsubscript{G}} garantisce che l'ambiente di sviluppo e produzione sia coerente e riproducibile. Questo approccio non solo assicura che le applicazioni funzionino correttamente in ogni fase del ciclo di vita, ma permette anche una gestione più efficiente delle dipendenze e delle configurazioni.\\
Di seguito un elenco delle best practice per la scrittura di file \href{https://7last.github.io/docs/pb/documentazione-interna/glossario\#docker}{Docker\textsubscript{G}}.
\begin{itemize}
	\item \textbf{Chiarezza e semplicità}: sono principi fondamentali nella scrittura di un Dockerfile. Questi principi aiutano a creare immagini \href{https://7last.github.io/docs/pb/documentazione-interna/glossario\#docker}{Docker\textsubscript{G}} che sono facili da comprendere, mantenere e utilizzare.
	\item \textbf{Versionamento}: è una pratica essenziale per mantenere il controllo sulle diverse versioni del software e delle dipendenze all'interno di un Dockerfile. Questa pratica aiuta a garantire la riproducibilità delle build, a facilitare il debug e a mantenere un ambiente di produzione stabile.
	\item \textbf{Sicurezza}: la sicurezza è una priorità cruciale nella costruzione di immagini \href{https://7last.github.io/docs/pb/documentazione-interna/glossario\#docker}{Docker\textsubscript{G}}. Un Dockerfile sicuro riduce la superficie di attacco, protegge dalle vulnerabilità e garantisce l'integrità dell'applicazione.
	\item \textbf{Efficienza}: nell'ottica di garantire tempi di risposta bassi, mantenendo comunque le prestazioni, i file \href{https://7last.github.io/docs/pb/documentazione-interna/glossario\#docker}{Docker\textsubscript{G}} devono avere un utilizzo corretto delle risorse e una gestione accurata dei container.
	\item \textbf{Gestione delle variabili d'ambiente}: è importante stabilire una pratica rigorosa che renda chiara la configurazione dell'applicazione all'interno del container. Le variabili d'ambiente sono fondamentali per la configurazione dinamica dell'applicazione, consentendo di adattare il comportamento dell'applicazione senza modificare direttamente il codice sorgente.
	\item \textbf{Monitoraggio}: riveste un'importanza fondamentale per garantire una corretta comprensione del comportamento dell'applicazione, individuare eventuali problemi e monitorare le prestazioni in tempo reale. Questo processo permette agli sviluppatori e agli \href{https://7last.github.io/docs/pb/documentazione-interna/glossario\#amministratore}{amministratori\textsubscript{G}} di sistema di ottenere una visione dettagliata delle attività dell'applicazione, della sua efficienza e di eventuali criticità.
	\item \textbf{Layering}: ogni istruzione nel Dockerfile crea un nuovo layer all'interno dell'immagine \href{https://7last.github.io/docs/pb/documentazione-interna/glossario\#docker}{Docker\textsubscript{G}}. La gestione efficiente di questi strati è cruciale poiché influisce direttamente sulle dimensioni dell'immagine \href{https://7last.github.io/docs/pb/documentazione-interna/glossario\#docker}{Docker\textsubscript{G}} finale e sulla velocità di build.
	\item \textbf{Riduzione delle dimensioni delle immagini}: la riduzione delle dimensioni delle immagini \href{https://7last.github.io/docs/pb/documentazione-interna/glossario\#docker}{Docker\textsubscript{G}} è un aspetto fondamentale per migliorare l'efficienza di deployment, ridurre i tempi di trasferimento e ottimizzare l'utilizzo delle risorse del sistema.
	\item \textbf{Documentazione}: fornire una documentazione esaustiva e chiara con un Dockerfile è essenziale per garantire la comprensione del suo funzionamento, facilitare la manutenzione e favorire la collaborazione tra i membri del team.
	\item \textbf{Testing}: testare un Dockerfile in modo completo è cruciale per garantire che l'immagine \href{https://7last.github.io/docs/pb/documentazione-interna/glossario\#docker}{Docker\textsubscript{G}} risultante sia configurata correttamente e funzioni come previsto.
\end{itemize}

\subsubsubsection{Strumenti}
\begin{itemize}
	\item \href{https://7last.github.io/docs/pb/documentazione-interna/glossario\#docker}{\textbf{Docker}\textsubscript{G}}: è una piattaforma open-source che permette di creare, distribuire e gestire applicazioni in ambienti di sviluppo e produzione utilizzando container leggeri e autonomi;
	\item \textbf{Visual Studio Code}: è un ambiente di sviluppo integrato (IDE) sviluppato da Microsoft e adottato dal gruppo per lo sviluppo del software. Offre funzionalità avanzate per la scrittura, la modifica e il debug del codice, con supporto per numerosi linguaggi di programmazione e framework.
\end{itemize}


\newpage

\section{Processi di supporto}
\subsection{Documentazione}
\subsubsection{Introduzione}
Il processo di documentazione è una componente fondamentale nella realizzazione e nel rilascio di un prodotto software,
poiché fornisce informazioni utili alle parti coinvolte e tiene traccia di tutte le attività relative al ciclo di vita del software,
comprese scelte e norme adottate dal gruppo durante lo svolgimento del progetto. In particolare, la documentazione è utile per:
\begin{itemize}
	\item permettere una comprensione profonda del prodotto e delle sue funzionalità;
	\item garantire uno standard di qualità all'interno dei processi produttivi.
\end{itemize}
Lo scopo della sezione è:
\begin{itemize}
	\item definire un insieme di regole e convenzioni per garantire la coerenza e la qualità della documentazione prodotta;
	\item creare template per ogni tipologia di documento così da garantire omogeneità.
\end{itemize}

\subsubsection{Versionamento}
La documentazione viene versionata allo stesso modo del codice sorgente, utilizzando la piattaforma \href{https://7last.github.io/docs/rtb/documentazione-interna/glossario\#github}{Github\textsubscript{G}}.
Ciò permette di tenere traccia delle modifiche apportate ai documenti e di realizzare procedure automatiche per l'individuazione degli errori grammaticali,
inserimento della "G" a pedice per i termini presenti nel \href{https://7last.github.io/docs/rtb/documentazione-interna/glossario\#glossario}{glossario\textsubscript{G}} e la pubblicazione nel sito del progetto.\\
Abbiamo stabilito di utilizzare il branch \texttt{main} per la pubblicazione dei documenti approvati dal \href{https://7last.github.io/docs/rtb/documentazione-interna/glossario\#verificatore}{verificatore\textsubscript{G}} e il branch \texttt{develop} per contenere
i file in formato LaTex da cui derivano i documenti finali.

\subsubsubsection{Branch policy}
Sono state definite le seguenti regole per la gestione dei branch:
\begin{itemize}
	\item \texttt{main}: viene aggiornato solamente dalla \href{https://7last.github.io/docs/rtb/documentazione-interna/glossario\#github}{\textit{Github\textsubscript{G}} Action} che si occupa di pubblicare i documenti approvati presenti nel branch \texttt{develop};
	\item \texttt{develop}: per la pubblicazione in questo branch è necessario ricevere un'approvazione da parte del \href{https://7last.github.io/docs/rtb/documentazione-interna/glossario\#verificatore}{verificatore\textsubscript{G}} tramite \textit{pull request}.
\end{itemize}

\subsubsection{Tipografia e sorgente documenti}
Per la redazione dei documenti abbiamo deciso di utilizzare il linguaggio di markup LaTeX, grazie alla sua flessibilità e possibilità di creare documenti di alta qualità tipografica.
Inoltre, abbiamo stabilito di suddividere ciascun documento in file differenti, in modo da consentire ai membri del team di lavorare contemporaneamente su sezioni diverse,
riducendo il numero di conflitti e semplificando la gestione del versionamento.

\subsubsection{Ciclo di vita}
Ciascun documento segue un ciclo di vita ben definito, che prevede le seguenti fasi:
\begin{enumerate}
	\item \textbf{Creazione branch} a partire da \texttt{develop}, utilizzando \href{https://www.atlassian.com/git/tutorials/comparing-workflows/gitflow-workflow}{\underline{Gitflow}} [Ultima consultazione 2024-05-13];
	\item \textbf{Scelta del template} da utilizzare per il documento, scegliendo tra \textit{documento}, \textit{verbale interno} o \textit{verbale esterno};
	\item \textbf{Redazione} del documento, seguendo le norme tipografiche e le convenzioni stabilite;
	\item \textbf{Commit} delle modifiche effettuate;
	\item \textbf{Apertura \textit{pull request}} per la revisione del documento, che dovrà essere approvata dal \href{https://7last.github.io/docs/rtb/documentazione-interna/glossario\#verificatore}{verificatore\textsubscript{G}} e dal \href{https://7last.github.io/docs/rtb/documentazione-interna/glossario\#responsabile}{responsabile\textsubscript{G}} per lo \href{https://7last.github.io/docs/rtb/documentazione-interna/glossario\#sprint}{sprint\textsubscript{G}} in corso;
	\item \textbf{Chiusura \textit{pull request}} e \textbf{merge} del branch;
	\item \textbf{Completamento task} su \href{https://7last.github.io/docs/rtb/documentazione-interna/glossario\#clickup}{\textit{ClickUp}\textsubscript{G}}.
\end{enumerate}

\subsubsection{Convenzioni nomenclatura e struttura di archiviazione}
Tutti i documenti prodotti devono seguire la seguente convenzione di nomenclatura ben definita, in modo da garantire una corretta organizzazione e archiviazione dei file
\begin{center}
	\texttt{nome\_fase/tipo\_documento/nome\_documento/nome\_documento.tex}
\end{center}
Dove:
\begin{itemize}
	\item \textbf{nome\_fase}: rappresenta la fase di progetto in cui il documento è stato prodotto, può assumere i seguenti valori:
	      \begin{itemize}
		      \item \textbf{1\_candidatura};
		      \item \textbf{2\_RTB};
		      \item \textbf{3\_PB};
		      \item \textbf{4\_CA}.
	      \end{itemize}
	\item \textbf{tipo\_documento}: rappresenta la tipologia di documento, può assumere i seguenti valori:
	      \begin{itemize}
		      \item \textbf{documentazione\_interna};
		      \item \textbf{documentazione\_esterna};
		      \item \textbf{verbali\_interni};
		      \item \textbf{verbali\_esterni}.
	      \end{itemize}
	\item \textbf{nome\_documento}: rappresenta il nome del documento, scritto in snake case.
\end{itemize}

\subsubsection{Procedure correlate alla redazione di documenti}
\subsubsubsection{I redattori}
Il redattore si occupa della stesura del documento, seguendo le norme tipografiche e le convenzioni stabilite, in modo da garantire la
chiarezza e la coerenza del testo. Dovrà rispettare le convenzioni sopracitate.

\subsubsubsection{Il resposabile}
Il \href{https://7last.github.io/docs/rtb/documentazione-interna/glossario\#responsabile}{responsabile\textsubscript{G}} coordina e gestisce le attività del gruppo. In particolare:
\begin{itemize}
	\item definire quali documenti devono essere prodotti e quali compiti devono essere svolti;
	\item gestisce le risorse e gli strumenti necessari per ultimare tali attività;
	\item approva la versione finale dei documenti;
	\item comunica con gli \href{https://7last.github.io/docs/rtb/documentazione-interna/glossario\#stakeholder}{stakeholder\textsubscript{G}}.
\end{itemize}

\subsubsubsection{L'amministratore}
L’\href{https://7last.github.io/docs/rtb/documentazione-interna/glossario\#amministratore}{amministratore\textsubscript{G}} è colui che si occupa della configurazione degli strumenti di supporto alla produzione del software; nel nostro caso,
inserisce le attività nell'ITS \href{https://7last.github.io/docs/rtb/documentazione-interna/glossario\#clickup}{ClickUp\textsubscript{G}} e le assegna ai membri del gruppo.

\subsubsubsection{I verificatori}
Il compito del \href{https://7last.github.io/docs/rtb/documentazione-interna/glossario\#verificatore}{verificatore\textsubscript{G}} è di controllare la correttezza e la qualità dei documenti prodotti, assicurandosi che siano conformi agli standard e alle norme stabilite.

\subsubsection{Struttura del documento}
Tutti i documenti produtti devono seguire una struttura ben definita, che prevede le componenti in seguito elencate.
\subsubsubsection{Prima pagina}
Nella prima pagina di ogni documento è presente un'intestazione contenente le seguenti informazioni:
\begin{itemize}
	\item titolo del documento;
	\item versione del documento;
	\item logo del gruppo;
	\item nome del gruppo.
\end{itemize}

\subsubsubsection{Registro delle modifiche}
La seconda pagina è dedicata al registro delle modifiche, rappresentato mediante una tabella. Il suo scopo è quello di tenere traccia di tutti i cambiamenti effettuati al documento,
in modo da garantire la rintracciabilità e la trasparenza delle modifiche. Ogni riga della tabella corrisponde a un cambiamento, e contiene le seguenti informazioni:
\begin{itemize}
	\item \textbf{versione} del documento, nel formato vX.Y, dove X rappresenta il numero di versione principale e Y il numero di versione secondaria;
	\item \textbf{data} di ultima modifica;
	\item nome del \textbf{redattore};
	\item nome del \href{https://7last.github.io/docs/rtb/documentazione-interna/glossario\#verificatore}{\textbf{verificatore}\textsubscript{G}};
	\item \textbf{descrizione} della modifica.
\end{itemize}

\subsubsubsection{Indice}
Ogni documento contiene un indice delle sezioni e delle sottosezioni presenti al suo interno, in modo da facilitare la consultazione e la navigazione.
Tutte le figure e tabelle presenti nei documenti sono visibili negli appositi indici e sono direttamente accessibili tramite link ipertestuali.

\subsubsubsection{Intestazione}
Ogni pagina del documento, ad eccezione della prima, include un'intestazione che presenta, da sinistra a destra, il logo del gruppo, il titolo del documento e la versione.

\subsubsubsection{Struttura dei verbali}
I verbali rappresentano un resoconto dettagliato degli incontri, tracciando gli argomenti discussi, le decisioni prese e le azioni da intraprendere.
Sono suddivisi in verbali interni ed esterni, a seconda che il meeting coinvolga solo i membri del team o persone esterne al gruppo.
La struttura è la medesima per entrambe le tipologie, ad eccezione dell'ultima pagina.
\begin{itemize}
	\item \textbf{prima pagina}:
	      \begin{itemize}
		      \item tipologia verbale (interno / esterno);
		      \item titolo del documento;
		      \item sottotitolo del documento, "Riunione interna settimanale" per i verbali interni e "Colloquio con \textit{Sync Lab}" per i verbali esterni;
		      \item data e versione del documento;
		      \item logo del gruppo
		      \item nome del gruppo.
	      \end{itemize}
	\item \textbf{Registro delle modifiche};
	\item \textbf{Dettagli sulla riunione}:
	      \begin{itemize}
		      \item sede della riunione (piattaforma utilizzata);
		      \item orario di inizio e fine;
		      \item partecipanti del gruppo (interni) con specificazione del ruolo;
		      \item partecipanti esterni.
	      \end{itemize}
	\item \textbf{Corpo del documento}
	      All'interno del corpo del documento sono presenti le seguenti sezioni:
	      \begin{itemize}
		      \item \textbf{ordine del giorno}: elenco di ciò che verrà discusso durante la riunione;
		      \item \textbf{verbale}: resoconto dettagliato dell'incontro, con una sottosezione per ciascun punto dell'ordine del giorno,
		            gli obiettivi prefissati e le decisioni prese;
	      \end{itemize}
	\item \textbf{Ultima pagina}: solo in caso di verbale esterno, l'ultima pagina contiene data e firma di un componente esterno, presente durante l'incontro.
\end{itemize}
Il template dei verbali è disponibile nella repository del gruppo, più precisamente al seguente \href{https://github.com/7Last/docs/tree/develop/0_template}{\underline{link}}. (Ultima consultazione 2024-05-20).
\newpage
\subsubsection{Norme tipografiche}
\textbf{Nomi dei file}\\ I nomi dei documenti devono essere coerenti con il loro contenuto,
scritti in snake case. Nel branch \texttt{develop} tutti i documenti non contengono la versione nel nome, al fine di evitare conflitti durante la modifica.
Nel momento in cui il documento viene approvato, esso viene spostato nel branch \texttt{main} e la versione viene aggiunta automaticamente al nome del file.\\
In particolare, la denominazione dei file nel branch \texttt{main} deve seguire la seguente convenzione:
\begin{itemize}
	\item \textbf{Verbali}: verbale\_esterno/interno\_YYYY\_MM\_DD\_vA.B;
	\item \href{https://7last.github.io/docs/rtb/documentazione-interna/glossario\#norme-di-progetto}{\textbf{Norme di Progetto}\textsubscript{G}}: norme\_di\_progetto\_vA.B;
	\item \href{https://7last.github.io/docs/rtb/documentazione-interna/glossario\#analisi-dei-requisiti}{\textbf{Analisi dei Requisiti}\textsubscript{G}}: analisi\_dei\_requisiti\_vA.B;
	\item \href{https://7last.github.io/docs/rtb/documentazione-interna/glossario\#piano-di-progetto}{\textbf{Piano di Progetto}\textsubscript{G}}: piano\_di\_progetto\_vA.B;
	\item \textbf{Lettera di Presentazione}: lettera\_di\_presentazione;
	\item \href{https://7last.github.io/docs/rtb/documentazione-interna/glossario\#glossario}{\textbf{Glossario}\textsubscript{G}}: \href{https://7last.github.io/docs/rtb/documentazione-interna/glossario\#glossario}{glossario\textsubscript{G}}\_vA.B.
\end{itemize}
\textbf{Stile del testo}
\begin{itemize}
	\item \textbf{Grassetto}:
	      \begin{itemize}
		      \item titoli di sezione;
		      \item termini importanti;
		      \item parole seguite da descrizione o elenchi puntati.
	      \end{itemize}
	\item \textbf{Corsivo}:
	      \begin{itemize}
		      \item nome del gruppo e dell'azienda \href{https://7last.github.io/docs/rtb/documentazione-interna/glossario\#proponente}{proponente\textsubscript{G}};
		      \item riferimenti a documenti esterni.
	      \end{itemize}
	\item \textbf{Maiuscolo}:
	      \begin{itemize}
		      \item acronimi;
		      \item iniziali dei nomi;
	      \end{itemize}
\end{itemize}
\newpage
\textbf{Regole sintattiche}:
\begin{itemize}
	\item negli elenchi ogni voce deve terminare con ";" mentre l'ultima ".";
	\item i numeri razionali vengono scritti mediante l'uso della virgola come separatore tra parte intera e parte decimale;
	\item le date devono seguire lo standard internazionale ISO 8601, rappresentando la data con YYYY-MM-DD (anno, mese, giorno).
	\item le sezioni di un documento devono essere numerate in modo gerarchico, seguendo la struttura X.Y, dove:
	      \begin{itemize}
		      \item X rappresenta il numero della sezione principale;
		      \item Y rappresenta il numero della sottosezione;
	      \end{itemize}
\end{itemize}

\subsubsection{Metriche}

\begin{table}[!h]
	\centering
	\begin{tabular}{|c|l|}
		\hline
		\textbf{Codice} & \textbf{Nome esteso}               					\\
		\hline
		\underline{\hyperlink{19M}{19M-IG}}     & Indice di Gulpease   			\\
		\underline{\hyperlink{20M}{20M-CO}}     & Correttezza Ortografica  		\\
		\hline
	\end{tabular}
	\caption{Metriche inerenti al processo di Documentazione}
\end{table}

\subsubsection{Strumenti}
Gli strumenti utilizzati per la redazione dei documenti sono:
\begin{itemize}
	\item \textbf{LaTeX}: impiegato nella stesura dei documenti;
	\item \href{https://7last.github.io/docs/rtb/documentazione-interna/glossario\#github}{\textbf{GitHub}\textsubscript{G}}: utilizzato per la gestione delle versioni e per la condivisione dei documenti;
	\item \textbf{Visual Studio Code}: utilizzato come editor di testo per la scrittura dei documenti.
\end{itemize}



\subsection{Accertamento della qualità}
\subsubsection{Introduzione}
L'accertamento della qualità consiste in un insieme di attività e processi volti a garantire che il software sviluppato soddisfi gli standard di qualitativi richiesti. Questo include l'adesione a requisiti concordati con la \href{https://7last.github.io/docs/rtb/documentazione-interna/glossario\#proponente}{proponente\textsubscript{G}}, normative ed è cruciale per assicurare che il prodotto finale sia affidabile e performante.

\subsubsection{Attività}
Il processo di accertamento della qualità comprende le seguenti attività:
\begin{enumerate}
	\item \textbf{sviluppo del \href{https://7last.github.io/docs/rtb/documentazione-interna/glossario\#piano-di-qualifica}{piano di qualifica\textsubscript{G}}}: vengono definiti gli standard di qualità e le metriche che il progetto deve raggiungere;
	\item \textbf{monitoraggio continuo della qualità}: consiste in attività sistematiche e pianificate per garantire che i processi utilizzati durante lo sviluppo del software siano adeguati e seguiti correttamente;
	\item \textbf{gestione delle configurazioni}: attraverso un sistema di gestione delle configurazioni ci si assicura che ciascuna di esse venga sottoposta ad opportuni controlli e valutazioni;
	\item \textbf{reporting}: in questa fase vengono generati report che descrivono lo stato della qualità del progetto, inclusi i risultati dei test, i difetti trovati e le azioni correttive intraprese;
	\item \textbf{miglioramento continuo}: vengono applicate le modifiche ai processi basate sui feedback e sulle analisi per migliorare continuamente la qualità del software e dei processi di sviluppo;
	\item \textbf{confronto con la \href{https://7last.github.io/docs/rtb/documentazione-interna/glossario\#proponente}{proponente\textsubscript{G}}}: viene instaurata una comunicazione con la \href{https://7last.github.io/docs/rtb/documentazione-interna/glossario\#proponente}{proponente\textsubscript{G}} per ricevere feedback costante sul lavoro svolto.
\end{enumerate}

\subsubsection{Piano di qualifica}
Il documento \href{https://7last.github.io/docs/rtb/documentazione-interna/glossario\#piano-di-qualifica}{\textit{Piano di Qualifica}\textsubscript{G}} è redatto per garantire che il software sviluppato rispetti gli standard di qualità richiesti e soddisfi le aspettative degli \href{https://7last.github.io/docs/rtb/documentazione-interna/glossario\#stakeholder}{stakeholder\textsubscript{G}}. Il suo utilizzo e il suo scopo si estendono a diverse aree critiche del progetto, dalle fasi iniziali di pianificazione fino alla consegna finale del prodotto.

\subsubsection{Ciclo di Deming}
Il \textit{ciclo di Deming}, è un ciclo a 4 stadi che consente di apportare determinate migliorie a specifici processi. È composto da:
\begin{itemize}
	\item \textbf{Plan}: in cui si definiscono le attività, scadenze, responsabilità e risorse per raggiungere gli obiettivi di miglioramento;
	\item \textbf{Do}: esecuzione delle attività definite dal punto precedente;
	\item \textbf{Check}: verifica l'esito delle azioni di miglioramento rispetto alle attese;
	\item \textbf{Act}: consolidare il buono e cercare modi per migliorare il resto.
\end{itemize}

È importante specificare che questo ciclo \textbf{non opera} sul prodotto, bensì sul \textbf{\textit{way of working}} del gruppo, per migliorarlo e renderlo più efficiente.

\subsubsection{Struttura e identificazioni metriche}
Ogni metrica presenta la seguente struttura:
\begin{itemize}
	\item \textbf{Metrica}:
	      codice identificativo nel formato:
	      \begin{center}
		      \textbf{[numero]M-[acronimo]}
	      \end{center}
	      Dove:
	      \begin{itemize}
		      \item \textbf{[numero]}: numero progressivo univoco per ogni metrica;
		      \item \textbf{M}: metrica;
		      \item \textbf{[acronimo]}: abbreviazione del nome della metrica.
	      \end{itemize}
	\item \textbf{Nome}: nome della metrica;
	\item \textbf{Valore accettabile}: valore minimo affinché la metrica sia considerabile soddisfacente e conforme agli obiettivi di qualità;
	\item \textbf{Valore ammissibile}: valore ottimale e ideale che dovrebbe essere raggiunto dalla metrica;
	\item \textbf{Valore ottimo}:
	\item \textbf{Descrizione}: breve descrizione della metrica adottata e delle sue funzionalità;
\end{itemize}

\subsubsection{Metriche}
\begin{table}[!h]
	\centering
	\begin{tabular}{ | c | l | }
		\hline
		\textbf{Codice}                      & \textbf{Nome esteso}            \\
		\hline
		\underline{\hyperlink{21M}{21M-FU}}  & Facilità di utilizzo            \\
        \underline{\hyperlink{22M}{22M-TA}}  & Tempo di apprendimento          \\
        \underline{\hyperlink{23M}{23M-TR}}  & Tempo di risposta               \\
        \underline{\hyperlink{24M}{24M-TE}}  & Tempo di elaborazione           \\
        \underline{\hyperlink{25M}{25M-QMS}} & Metriche di qualità soddisfatte \\
		\hline
	\end{tabular}
	\caption{Metriche relative all'Accertamento della qualità}
\end{table}

\subsection{Verifica}
\subsubsection{Introduzione}
La verifica nel ciclo di vita del software garantisce l'efficienza e la correttezza delle attività, assicurando che i prodotti soddisfino i requisiti. I \href{https://7last.github.io/docs/rtb/documentazione-interna/glossario\#verificatore}{verificatori\textsubscript{G}} applicano tecniche di test seguendo procedure definite all'interno del \href{https://7last.github.io/docs/rtb/documentazione-interna/glossario\#piano-di-qualifica}{Piano di Qualifica\textsubscript{G}}, il quale traccia il percorso di verifica, stabilendo obiettivi e criteri di accettazione. Tutti i processi attivi vengono sottoposti a verifica dopo aver raggiunto un buon grado di completamento o a seguito di modifiche dello stato del processo stesso.

\subsubsection{Verifica dei documenti} \label{verifica_dei_documenti}
Nell'ambito della documentazione, la verifica è un'attività cruciale per garantire la correttezza e l'accuratezza dei contenuti. Questo processo coinvolge una serie di controlli sistematici e revisioni finalizzate a garantire che la documentazione sia accurata, completa e conforme agli standard e ai requisiti specificati.
La verifica dei documenti assicura l'accuratezza, identificando e correggendo errori, omissioni e incongruenze. Questo riduce il rischio di malintesi o errori di implementazione, garantendo che tutte le informazioni tecniche e di progetto siano corrette e precise.
Inoltre, la verifica garantisce che la documentazione rispetti gli standard aziendali, di settore e normativi, mantenendo la coerenza e la qualità attraverso tutti i documenti del progetto. Documenti verificati e accurati migliorano la comunicazione tra i membri del team e con gli \href{https://7last.github.io/docs/rtb/documentazione-interna/glossario\#stakeholder}{stakeholder\textsubscript{G}} esterni, facilitando la comprensione e la collaborazione e riducendo i tempi di chiarimento e discussione. Essa si suddivide in:
\begin{itemize}
	\item \textbf{correttezza tecnica}: vengono controllate le informazioni tecniche e i contenuti del documento, verificando che siano corretti e precisi;
	\item \textbf{conformità agli standard}: verifica che il documento segua le linee guida e gli standard stabiliti per la formattazione, la struttura e lo stile;
	\item \textbf{verifica della completezza}: si effettua un controllo per assicurare che tutti i requisiti documentali siano stati soddisfatti e che  tutte le sezioni e le informazioni necessarie siano presenti;
	\item \textbf{revisione linguistica}: viene effettuata una correzione di eventuali errori grammaticali, ortografici e di punteggiatura con l'ausilio di controlli automatici;
	\item \textbf{coerenza}: ci si assicura che i termini e le definizioni siano utilizzati in modo uniforme e coerente in tutto il documento.
\end{itemize}

\subsubsubsection{Processo di revisione}
Al momento di completamento delle modifiche in un documento il redattore è tenuto a creare una pull request dal branch in cui ha effettuato le modifiche verso il branch \texttt{develop}, richiedendo una revisione da parte del \href{https://7last.github.io/docs/rtb/documentazione-interna/glossario\#verificatore}{verificatore\textsubscript{G}} e del \href{https://7last.github.io/docs/rtb/documentazione-interna/glossario\#responsabile}{responsabile\textsubscript{G}}; in caso di approvazione verrà effettuato il merge e la issue creata su \href{https://7last.github.io/docs/rtb/documentazione-interna/glossario\#clickup}{ClickUp\textsubscript{G}} verrà spostata da \textit{in progres} a \textit{done}.

\subsubsection{Analisi}
L'analisi è un processo essenziale che mira a garantire la qualità e l'affidabilità del prodotto finale. Questa fase critica comporta un esame dettagliato e sistematico di tutte le componenti del progetto per identificare e risolvere eventuali problemi prima della consegna. Si suddivide in \textbf{analisi dinamica} e \textbf{analisi statica}.

\subsubsubsection{Analisi dinamica}
L'analisi dinamica riguarda l'esame del comportamento del software durante la sua esecuzione. Questo tipo di analisi viene eseguito a tempo di esecuzione e si concentra sull'osservazione di come il software interagisce con il sistema operativo, le risorse di sistema e altre applicazioni. L'analisi dinamica è essenziale per valutare l'efficienza, le prestazioni e la sicurezza del software in condizioni operative reali; viene effettuata tramite l'utilizzo delle seguenti tipologie di test:
\begin{itemize}
	\item test di unità;
	\item test di integrazione;
	\item test di sistema;
	\item test di regressione;
	\item test di accettazione.
\end{itemize}

I benefici dell'analisi dinamica includono:
\begin{itemize}
	\item \textbf{individuazione di errori e difetti}: consente di identificare problemi e difetti nel software durante l'esecuzione, riducendo i rischi di malfunzionamenti e guasti;
	\item \textbf{ottimizzazione delle prestazioni}: permette di valutare le prestazioni del software e di identificare possibili aree di miglioramento;
	\item \textbf{convalida delle funzionalità}: assicura che il software funzioni correttamente e soddisfi i requisiti operativi attraverso i test elencati in precedenza.
\end{itemize}

\hypertarget{testing}{\subsubsection{Testing}}
Il testing è una fase fondamentale dello sviluppo del software, il cui scopo primario è garantire che il prodotto finale sia di alta qualità, soddisfi i requisiti specificati e funzioni correttamente in tutte le condizioni previste. Questo processo si articola in una serie di attività sistematiche e metodiche per identificare difetti, errori e incongruenze nel software, migliorandone così l'affidabilità e le prestazioni complessive.

Per ogni test è necessario definire i seguenti aspetti:
\begin{itemize}
	\item \textbf{identificazione del test}: un identificatore univoco per il caso di test, che faciliti il tracciamento e la referenza durante il processo di testing;
	\item \textbf{descrizione del test}: una descrizione chiara e concisa del caso di test, che spieghi la funzionalità specifica o il requisito che viene testato;
	\item \textbf{input del test}: dettagli sui dati di input necessari per eseguire il test, inclusi eventuali prerequisiti o condizioni iniziali;
	\item \textbf{output atteso}: descrizione degli output o dei risultati attesi del test, inclusi valori attesi, comportamenti del sistema o condizioni post-esecuzione;
	\item \textbf{stato del test}: indicazione dello stato del test, che può essere "superato", "non superato" o "non implementato";
\end{itemize}

\subsubsubsection{Test di unità}
I test di unità consentono di valutare il funzionamento delle singole unità o componenti del codice in modo isolato, senza dipendere dalle altre parti del sistema. Questo approccio mira a garantire che ogni porzione di software, anche la più piccola, operi correttamente e coerentemente con le specifiche definite, indipendentemente dalle interazioni con altre unità.
L'obiettivo principale dei test di unità è di assicurare che ogni unità funzioni come previsto, restituendo risultati attesi per determinate condizioni di input. Ciò significa che, per ciascuna unità, vengono definiti uno o più scenari di test che includono input specifici e i risultati attesi per tali input. Durante l'esecuzione dei test, il comportamento effettivo dell'unità viene confrontato con i risultati attesi, e qualsiasi discrepanza tra i due indica un potenziale difetto nel codice.
La pratica dei test di unità offre diversi vantaggi significativi. Innanzitutto, permette di individuare e risolvere tempestivamente eventuali difetti nel codice, poiché consente di identificare problemi già durante le prime fasi dello sviluppo, quando sono più facili e meno costosi da correggere. Inoltre, l'esecuzione automatica dei test di unità consente di garantire una verifica continua del codice, permettendo agli sviluppatori di individuare rapidamente eventuali regressioni o effetti collaterali indesiderati derivanti da modifiche al codice.
Inoltre, i test di unità forniscono una documentazione vivente del comportamento del software, in quanto ciascun test rappresenta un caso specifico di utilizzo dell'unità. Questo non solo facilita la comprensione del codice da parte di altri sviluppatori, ma può anche servire come una sorta di contratto che definisce il comportamento atteso dell'unità nel tempo.

\subsubsubsection{Test di integrazione}
I test di integrazione sono una fase nel processo di testing del software durante la quale vengono verificati e validati i collegamenti e le interazioni tra le diverse unità o componenti del sistema. L'obiettivo principale dei test di integrazione è accertarsi che le singole parti del software, già testate individualmente durante i test di unità, funzionino correttamente quando integrate insieme come un sistema completo. In pratica, i test di integrazione esaminano come le varie unità o componenti interagiscono tra loro e se cooperano correttamente per fornire le funzionalità desiderate del software. Questi test possono rivelare problemi come:
\begin{itemize}
	\item \textbf{incompatibilità} tra le interfacce delle diverse unità;
	\item \textbf{errori di comunicazione} tra le componenti;
	\item \textbf{mancata sincronizzazione o coerenza} nei dati scambiati tra le unità;
\end{itemize}

\subsubsubsection{Test di sistema}
Durante questa tipologia di test l'intero sistema viene valutato e verificato per garantire che soddisfi i requisiti funzionali e non funzionali specificati. Vengono eseguiti dopo i test di unità e di integrazione e sono progettati per valutare il sistema nel suo complesso, testando le sue funzionalità e le sue prestazioni rispetto agli obiettivi definiti. In particolare, permettono di:
\begin{itemize}
	\item \textbf{verificare} che il sistema \textbf{soddisfi i requisiti specificati};
	\item \textbf{valutare le prestazioni} del sistema;
	\item \textbf{verificare la sicurezza} e la \textbf{robustezza} del sistema;
	\item \textbf{verificare la compatibilità} del sistema con l'ambiente di esecuzione;
	\item \textbf{verificare la facilità d'uso} e l'\textbf{usabilità} del sistema.
\end{itemize}

\subsubsubsection{Test di regressione}
I test di regressione assicurano che le modifiche apportate al codice non abbiano introdotto nuovi difetti o intaccato funzionalità esistenti nel software. Questi test vengono eseguiti dopo che sono state apportate modifiche al software, come nuove funzionalità, correzioni di bug o aggiornamenti del sistema, per assicurare il corretto funzionamento anche dopo tali modifiche.

\subsubsubsection{Test di accettazione}
I test di accettazione sono una fase nel processo di sviluppo del software durante la quale il prodotto viene testato per valutare se soddisfa i requisiti e le aspettative degli utenti finali, dei clienti o degli \href{https://7last.github.io/docs/rtb/documentazione-interna/glossario\#stakeholder}{stakeholder\textsubscript{G}}. Questi test vengono eseguiti dopo che il software è stato completamente sviluppato e prima del suo rilascio ufficiale, consentendo agli utenti finali di valutarlo in un ambiente simile a quello di produzione e fornire feedback sulle sue funzionalità, l'usabilità e la conformità ai requisiti.

\subsubsubsection{Sequenza delle fasi di test}
I test vengono eseguiti in sequenza, partendo da quelli di unità e procedendo verso quelli di integrazione, di sistema, di regressione e di accettazione. Questo approccio graduale consente di identificare e risolvere i difetti in modo sistematico e strutturato, garantendo che il software funzioni correttamente e soddisfi i requisiti specificati.

\subsubsubsection{Codici dei test}
Ciascun test deve essere identificato da un codice univoco nel seguente formato:
\begin{center}
	\textbf{[numero\_test]-[tipologia]}
\end{center}
Dove:
\begin{itemize}
	\item \textbf{[numero\_test]}: rappresenta un numero, in ordine crescente, associato al test, univoco all'interno della tipologia;
	\item \textbf{[tipologia]}: indica l'appartenenza del test, può assumere i seguenti valori:
	      \begin{itemize}
		      \item \textbf{U}: test di unità;
		      \item \textbf{I}: test di integrazione;
		      \item \textbf{S}: test di sistema;
		      \item \textbf{R}: test di regressione;
		      \item \textbf{A}: test di accettazione.
	      \end{itemize}

\end{itemize}

\subsubsubsection{Stato dei test}
A ciascun test è associato uno stato che indica l'esito della sua esecuzione:
\begin{itemize}
	\item \textbf{S}: il test è stato superato;
	\item \textbf{NS}: il test non è stato superato;
	\item \textbf{NI}: il test non è stato implementato.
\end{itemize}
Questi risultati saranno riportati nel documento \textit{"\href{https://7last.github.io/docs/rtb/documentazione-interna/glossario\#piano-di-qualifica}{Piano di Qualifica\textsubscript{G}}"}, in particolare nella sezione \textit{"Metodologie di Testing"}.

\subsubsubsection{Analisi statica}
L'analisi statica è il processo mediante il quale si esamina il codice sorgente di un programma senza eseguirlo. Questo tipo di analisi viene effettuato utilizzando strumenti automatizzati che analizzano il codice per individuare possibili errori, violazioni delle best practice di codifica e potenziali vulnerabilità di sicurezza. L'analisi statica si concentra su aspetti come la sintassi, la semantica e le strutture di controllo del codice. I benefici di questo tipo di analisi sono molteplici, tra cui:
\begin{itemize}
	\item \textbf{individuazione precoce degli errori}: consente di rilevare errori e difetti nelle prime fasi dello sviluppo, riducendo i costi e i tempi necessari per le correzioni;
	\item \textbf{miglioramento della qualità del codice}: promuove l'adozione di standard di codifica e best practice, migliorando la manutenibilità e la leggibilità del codice;
	\item \textbf{prevenzione di vulnerabilità di sicurezza}: aiuta a individuare vulnerabilità di sicurezza, come buffer overflow e injection flaws, prima che il software venga eseguito.
\end{itemize}

Nell'ambito dell'analisi statica del software, l'inspection e il walkthrough sono tecniche di revisione del codice e della documentazione, ma presentano differenze significative in termini di obiettivi, formalità e procedure.

\subsubsubsubsection{Inspection}
L'inspection è una tecnica formale e strutturata di revisione del codice. Le sue principali caratteristiche sono:
\begin{itemize}
	\item \textbf{formalità}: l'inspection segue un processo rigoroso con ruoli definiti (moderatore, autore, lettore, ispettore, e registratore) e fasi ben precise (pianificazione, overview, preparazione, meeting di inspection, rework, e follow-up);
	\item \textbf{obiettivi}: l'obiettivo principale è identificare difetti nel codice o nella documentazione. Si concentra su aspetti come errori di logica, violazioni degli standard di codifica, problemi di performance e sicurezza;
	\item \textbf{documentazione}: viene prodotta una documentazione dettagliata delle osservazioni e dei difetti riscontrati, che serve come base per la correzione e per tracciare le azioni successive;
	\item \textbf{ruoli e responsabilità}: ogni partecipante ha un ruolo specifico e le attività sono strettamente coordinate. Il moderatore gestisce il processo, l'autore presenta il lavoro, gli ispettori individuano i difetti, il lettore guida la revisione del materiale, e il registratore annota i difetti trovati;
	\item \textbf{preparazione} I partecipanti devono esaminare il materiale da rivedere in anticipo e preparare una lista di potenziali problemi.
\end{itemize}

\subsubsubsubsection{walkthrough}
Il walkthrough è una tecnica meno formale e più flessibile rispetto all'inspection. Le sue principali caratteristiche sono:
\begin{itemize}
	\item \textbf{informalità}: il walkthrough è meno strutturato e può essere condotto in modo informale. Non richiede ruoli rigidamente definiti o una procedura rigorosa;
	\item \textbf{obiettivi}: l'obiettivo principale è comprendere il codice o la documentazione, discutere possibili miglioramenti e condividere conoscenze tra i membri del team. Anche se l'individuazione dei difetti è uno degli scopi, non è l'unico focus;
	\item \textbf{documentazione}: la documentazione delle osservazioni è meno formale e dettagliata rispetto all'inspection. Può non essere sempre prodotta una documentazione completa delle discussioni e dei difetti riscontrati;
	\item \textbf{ruoli e responsabilità}: non ci sono ruoli formali assegnati. Solitamente, l'autore del codice o della documentazione guida la revisione e i partecipanti contribuiscono con commenti e suggerimenti;
	\item \textbf{preparazione}: la preparazione è meno intensiva. Non è necessario che tutti i partecipanti esaminino il materiale in anticipo, anche se può essere utile.
\end{itemize}

\subsubsection{Metriche}

\begin{table}[!h]
	\centering
	\begin{tabular}{|c|l|}
		\hline
		\textbf{Codice} & \textbf{Nome esteso}               					\\
		\hline
		\underline{\hyperlink{26M}{26M-CC}}     & Code Coverage   				\\
		\underline{\hyperlink{27M}{27M-BC}}     & Branch Coverage   			\\
		\underline{\hyperlink{28M}{28M-SC}}     & Statement Coverage   			\\
		\underline{\hyperlink{29M}{29M-FD}}     & Failure Density   			\\
		\underline{\hyperlink{30M}{30M-PTCP}}   & Passed Test Case Percentage   \\
		\hline
	\end{tabular}
	\caption{Metriche inerenti al processo di Verifica}
\end{table}



\subsubsection{Strumenti}

\begin{itemize}
	\item \href{http://aspell.net/}{Aspell} per il controllo ortografico;
	\item \href{https://docs.python.org/3/library/unittest.html}{unittest} libreria di \href{https://7last.github.io/docs/rtb/documentazione-interna/glossario\#python}{Python\textsubscript{G}} per lo sviluppo di test.
\end{itemize}



\subsection{Validazione}
\subsubsection{Introduzione}
La validazione rappresenta un momento critico nel ciclo di sviluppo, in quanto sottopone il software a una serie di controlli dettagliati per assicurare la sua conformità ai requisiti stabiliti e la sua idoneità all'utilizzo da parte degli utenti finali. Questo processo non è solo una verifica formale, ma una fase cruciale che assicura che il software sia costruito in modo tale da rispondere pienamente alle esigenze e alle aspettative degli utenti, nonché agli obiettivi del progetto.

\subsubsection{Procedura di validazione}
In questo processo copre un ruolo fondamentale il test di accettazione che mira a garantire la validazione
del prodotto. Infatti i diversi test elencati nella sezione \underline{\hyperlink{testing}{testing}} costituiscono un input
per la validazione. Essi dovranno verificare:
\begin{itemize}
	\item l'implementazione di tutti i casi d'uso;
	\item la conformità del prodotto ai requisiti obbligatori;
	\item il soddisfacimento di altri requisiti concordati con il \href{https://7last.github.io/docs/rtb/documentazione-interna/glossario\#committente}{committente\textsubscript{G}}.
\end{itemize}

\subsection{Gestione della configurazione}
\subsubsection{Introduzione}
La gestione della configurazione è il processo di identificazione, controllo e coordinamento dei componenti software e delle risorse associate durante tutto il ciclo di vita del prodotto. Questo processo assicura che il software e i suoi artefatti correlati siano gestiti in modo coerente e controllato, consentendo agli sviluppatori di tracciare le modifiche, gestire le versioni e garantire l'integrità e la coerenza del sistema nel tempo.
\newpage
\subsubsection{Versionamento}
La convenzione di versionamento adottata è nel formato X.Y dove:
\begin{itemize}
	\item \textbf{X}: rappresenta il completamento in vista di una delle fasi del progetto e dunque viene incrementato al raggiungimento di \href{https://7last.github.io/docs/rtb/documentazione-interna/glossario\#requirements-and-technology-baseline}{RTB\textsubscript{G}}, \href{https://7last.github.io/docs/rtb/documentazione-interna/glossario\#product-baseline}{PB\textsubscript{G}} ed eventuale \href{https://7last.github.io/docs/rtb/documentazione-interna/glossario\#customer-acceptance}{CA\textsubscript{G}}.
	\item \textbf{Y}: rappresenta una versione intermedia e viene incrementata ad ogni modifica significativa del documento.
\end{itemize}

\subsubsection{Repository}
Il team utilizza due repository:
\begin{itemize}
	\item \href{https://github.com/7Last/docs.git}{\underline{docs}}: contenente la documentazione prodotta;
	\item \href{https://github.com/7Last/SyncCity}{\href{https://7last.github.io/docs/rtb/documentazione-interna/glossario\#synccity}{\underline{SyncCity}\textsubscript{G}}}: contenente il codice del progetto.
\end{itemize}

\subsubsubsection{Struttura repository}
Il repository inerente alla documentazione è così organizzato:
\begin{enumerate}
	\item \textbf{Candidatura}:
	      \begin{itemize}
		      \item \textbf{verbali\_esterni}: al suo interno sono presenti i verbali delle riunioni avute con i membri di \textit{Sync Lab S.r.l.};
		      \item \textbf{verbali\_interni}: contenente i verbali delle riunioni interne svolte;
		      \item \textbf{lettera\_di\_presentazione}: documento di presentazione per la candidatura al \href{https://7last.github.io/docs/rtb/documentazione-interna/glossario\#capitolato}{capitolato\textsubscript{G}} scelto;
		      \item \textbf{preventivo\_costi\_assunzione\_impegni}: documento in cui vengono specificati i costi previsti, il totale delle ore per persona e gli impegni assunti;
		      \item \textbf{valutazione\_dei\_capitolati}: contenente il parere personale riguardo i capitolati offerti dalle varie aziende.
	      \end{itemize}
	\item \href{https://7last.github.io/docs/rtb/documentazione-interna/glossario\#requirements-and-technology-baseline}{\textbf{RTB}\textsubscript{G}}:
	      \begin{itemize}
		      \item \textbf{documentazione\_esterna}: contenente i seguenti documenti:
		            \begin{itemize}
			            \item \textbf{analisi\_dei\_requisiti};
			            \item \textbf{piano\_di\_progetto};
			            \item \textbf{piano\_di\_qualifica};
		            \end{itemize}
		      \item \textbf{documentazione\_interna}: contenente i seguenti documenti:
		            \begin{itemize}
			            \item \href{https://7last.github.io/docs/rtb/documentazione-interna/glossario\#glossario}{\textbf{glossario}\textsubscript{G}};
			            \item \textbf{norme\_di\_progetto}.
		            \end{itemize}
		      \item \textbf{verbali\_esterni}: al suo interno sono presenti i verbali delle riunioni avute con i membri di \textit{Sync Lab S.r.l.};
		      \item \textbf{verbali\_interni}: contenente i verbali delle riunioni interne svolte.
	      \end{itemize}
	\item \href{https://7last.github.io/docs/rtb/documentazione-interna/glossario\#product-baseline}{\textbf{PB}\textsubscript{G}} %todo, da implementare
\end{enumerate}

\subsubsection{Sincronizzazione e branching}

\subsubsubsection{Documentazione}	\label{nomenclatura}
L'approccio adottato per la redazione della documentazione segue il workflow noto come \textit{Gitflow}. Questo workflow è un processo strutturato e collaborativo che consente ai membri del team di lavorare in modo efficace alla creazione, revisione e integrazione di documenti all'interno di un progetto. Questo approccio garantisce una gestione ordinata e controllata della redazione dei documenti, consentendo una migliore organizzazione e una maggiore qualità del lavoro finale.
\begin{flushleft}
	\textbf{Nomenclatura dei branch per le attività di redazione e/o modifica di documenti}: \label{convenzioni_nomenclatura}
\end{flushleft}
\begin{itemize}
	\item il nome del nuovo branch deve riportare il titolo del documento da redarre o modificare;
	\item il nome dei verbali deve presentare anche la data della riunione:
	      \begin{itemize}
		      \item \textit{verbale\_interno\_yy\_mm\_dd} (es. verbale\_interno\_24\_03\_05);
	      \end{itemize}
\end{itemize}

\subsubsubsection{Sviluppo}
\textit{7Last} utilizza \href{https://www.atlassian.com/it/git/tutorials/comparing-workflows/gitflow-workflow}{\underline{{Gitflow}}} come flusso di lavoro.
\begin{flushleft}
	Flusso di lavoro Gitflow:
\end{flushleft}
\begin{enumerate}
    \item \textbf{\texttt{main} branch}
    \begin{itemize}
        \item branch principale e stabile;
        \item contiene il codice in produzione;
        \item ogni commit rappresenta una versione rilasciata.
    \end{itemize}

    \item \textbf{\texttt{develop} branch}
    \begin{itemize}
        \item branch per lo sviluppo integrato;
        \item riceve i merge dai feature branches;
        \item rappresenta lo stato intermedio prima del rilascio.
    \end{itemize}

    \item \textbf{feature branch}
    \begin{itemize}
        \item derivano dal branch \texttt{develop};
        \item utilizzati per sviluppare nuove funzionalità;
        \item dopo il completamento, viene fatto il merge nel branch \texttt{develop}.
    \end{itemize}

    \item \textbf{release branch}
    \begin{itemize}
        \item derivano dal branch \texttt{develop};
        \item preparano una nuova versione per il rilascio;
        \item permettono di effettuare correzioni di bug e preparare la documentazione;
        \item dopo il completamento, viene fatto il merge sia nel branch \texttt{main} che nel branch \texttt{develop}.
    \end{itemize}

    \item \textbf{hotfix branch}
    \begin{itemize}
        \item derivano dal branch master;
        \item utilizzati per correggere rapidamente bug critici in produzione;
        \item dopo il completamento, viene fatto il merge sia nel branch \texttt{main} che nel branch \texttt{develop}.
    \end{itemize}
\end{enumerate}

\subsubsection{Consegna e rilascio}
Il processo di consegna e rilascio secondo lo standard \textit{ISO/IEC 12207:1995} è una componente essenziale del ciclo di vita del software, che si concentra sulla gestione delle versioni del software e sulla distribuzione dei prodotti software agli utenti finali. Questo processo assicura che le versioni del software siano rilasciate in modo controllato, pianificato e che soddisfino i requisiti di qualità e sicurezza. Attraverso una serie di attività ben definite, il processo di gestione del rilascio garantisce che il software sia pronto per l'uso, minimizzando i rischi di problemi post-rilascio e massimizzando l'efficienza operativa. La creazione della prima release deve avvenire in concomitanza con la baseline \href{https://7last.github.io/docs/rtb/documentazione-interna/glossario\#requirements-and-technology-baseline}{RTB\textsubscript{G}} (\href{https://7last.github.io/docs/rtb/documentazione-interna/glossario\#requirements-and-technology-baseline}{Requirements and Technology Baseline\textsubscript{G}}), mentre la creazione delle successive release deve avvenire in concomitanza con la baseline \href{https://7last.github.io/docs/rtb/documentazione-interna/glossario\#product-baseline}{PB\textsubscript{G}} (\href{https://7last.github.io/docs/rtb/documentazione-interna/glossario\#product-baseline}{Product Baseline\textsubscript{G}}) e, se prevista, \href{https://7last.github.io/docs/rtb/documentazione-interna/glossario\#customer-acceptance}{CA\textsubscript{G}} (\href{https://7last.github.io/docs/rtb/documentazione-interna/glossario\#customer-acceptance}{Customer Acceptance\textsubscript{G}}).

\subsubsection{Sito del gruppo}
Il gruppo ha sviluppato un sito disponibile al link: \url{https://7last.github.io/}; che ha lo scopo di facilitare la consultazione dei documenti redatti fino ad ora. Tale pagina web viene aggiornata automaticamente per rispecchiare lo stato dei documenti presenti all'interno del branch \texttt{main}. Inoltre facilita la consultazione del \href{https://7last.github.io/docs/rtb/documentazione-interna/glossario\#glossario}{glossario\textsubscript{G}} grazie a dei link, presenti nelle parole distinte dalla lettera G al pedice, diretti alla definizione del termine.

\subsubsection{Strumenti}

\begin{itemize}
	\item \href{https://7last.github.io/docs/rtb/documentazione-interna/glossario\#github}{\textbf{Github}\textsubscript{G}}: utilizzato per la gestione delle versioni e per la condivisione dei documenti;
	\item \href{https://7last.github.io/docs/rtb/documentazione-interna/glossario\#clickup}{\textbf{ClickUp}\textsubscript{G}}: è una piattaforma di gestione del lavoro all-in-one progettata per pianificare, organizzare e collaborare su progetti e attività.
\end{itemize}


\subsection{Analisi congiunta}
\subsubsection{Introduzione}
L'analisi congiunta è un'attività collaborativa cruciale in cui i recensori e i recensiti si incontrano per esaminare e valutare diversi aspetti del progetto.  Nel nostro caso i recensori sono costituiti da \href{https://7last.github.io/docs/rtb/documentazione-interna/glossario\#proponente}{proponente\textsubscript{G}}, \href{https://7last.github.io/docs/rtb/documentazione-interna/glossario\#committente}{committente\textsubscript{G}} e \href{https://7last.github.io/docs/rtb/documentazione-interna/glossario\#stakeholder}{stakeholder\textsubscript{G}}; mentre i recensiti sono rappresentati dal gruppo \textit{7Last}. Lo scopo principale di questa revisione congiunta è garantire che tutte le parti coinvolte abbiano una comprensione comune dello stato attuale del progetto e dei passi successivi necessari per raggiungere gli obiettivi fissati.

\subsubsection{Realizzazione del processo}
Il processo comprende i seguenti impegni:
\subsubsubsection{Revisioni periodiche}
In corrispondenza delle \href{https://7last.github.io/docs/rtb/documentazione-interna/glossario\#milestone}{milestone\textsubscript{G}} stabilite saranno effettuate delle revisioni periodiche, come riportato nel
documento \href{https://7last.github.io/docs/rtb/documentazione-interna/glossario\#piano-di-progetto}{\textit{Piano di Progetto}\textsubscript{G}}.

\subsubsubsection{Stato Avanzamento Lavori}
Al termine di ogni \href{https://7last.github.io/docs/rtb/documentazione-interna/glossario\#sprint}{sprint\textsubscript{G}} viene svolta una revisione \href{https://7last.github.io/docs/rtb/documentazione-interna/glossario\#stato-avanzamento-lavori}{SAL\textsubscript{G}}
(\href{https://7last.github.io/docs/rtb/documentazione-interna/glossario\#stato-avanzamento-lavori}{Stato Avanzamento Lavori\textsubscript{G}}) tra il team e la \href{https://7last.github.io/docs/rtb/documentazione-interna/glossario\#proponente}{proponente\textsubscript{G}}. Questa revisione ha lo scopo di valutare il lavoro svolto durante lo \href{https://7last.github.io/docs/rtb/documentazione-interna/glossario\#sprint}{sprint\textsubscript{G}} precedente, per verificare che gli obiettivi prefissati siano stati correttamente raggiunti secondo le scadenze prefissate.
Inoltre, durante questo incontro, si pianificano le attività per lo \href{https://7last.github.io/docs/rtb/documentazione-interna/glossario\#sprint}{sprint\textsubscript{G}} successivo.

\subsubsubsection{Revisioni straordinarie}
Se uno qualsiasi dei soggetti coinvolti tra gli \href{https://7last.github.io/docs/rtb/documentazione-interna/glossario\#stakeholder}{stakeholder\textsubscript{G}} lo ritenga opportuno, è possibile istituire una revisione straordinaria per esaminare attentamente lo stato di avanzamento dei lavori. Durante questa revisione, si discutono eventuali problematiche emerse e le relative soluzioni adottabili.

\subsubsubsection{Risorse per le revisioni}
Le risorse coinvolte nelle revisioni possono essere differenti, ad esempio: strumenti hardware, software, documentazione... è fondamentale che tali risorse siano discusse e concordate tra tutte le parti coinvolte.

\subsubsubsection{Organizzazione degli incontri}
Pochi giorni prima della riunione con la \href{https://7last.github.io/docs/rtb/documentazione-interna/glossario\#proponente}{proponente\textsubscript{G}} \textit{7Last} si impegna a consegnare un report con tutte le modifiche effettuate e i campi in cui sono state applicate.
In ciascuna revisione rimane la necessità di concordare i seguenti elementi:
\begin{itemize}
	\item agenda della riunione;
	\item prodotti software risultati dall'attività e relative problematiche;
\end{itemize}

\subsubsubsection{Documenti prodotti e decisioni approvate}
Tramite i \textit{Verbali Esterni} vengono documentati i risultati delle revisioni, compresi i problemi individuati, le soluzioni proposte e le azioni correttive da intraprendere. Questi documenti vengono distribuiti a tutte le parti coinvolte per garantire una comprensione comune e una chiara comunicazione. La parte recensente comunicherà alla parte recensita la veridicità di quanto riportato, approvando o disapprovando i documenti citati.

\subsection{Risoluzione dei problemi}
\subsubsection{Introduzione}
La risoluzione dei problemi è un processo che mira a identificare, analizzare e risolvere le varie problematiche che possono emergere durante lo sviluppo. Si riferisce a un insieme di tecniche e metodologie utilizzate per affrontare e risolvere le difficoltà tecniche, di design, di implementazione o di gestione che possono sorgere durante lo sviluppo del software. Questo processo può includere:
\begin{itemize}
	\item \textbf{identificazione del problema}: consiste nel riconoscere e definire chiaramente il problema;
	\item \textbf{analisi del problema}: nella quale si esamina il problema per capire le cause sottostanti e l'impatto;
	\item \textbf{generazione di soluzioni}: vengono proposte diverse possibili soluzioni al problema;
	\item \textbf{valutazione delle soluzioni}: vengono analizzate le soluzioni proposte per determinarne la fattibilità, l'efficacia e l'efficienza;
	\item \textbf{implementazione della soluzione}: la soluzione adottata viene implementata e testata per verificare che risolva il problema in modo efficace.
\end{itemize}

\subsubsection{Gestione dei rischi}
All'interno del documento \href{https://7last.github.io/docs/rtb/documentazione-interna/glossario\#piano-di-progetto}{\textit{Piano di Progetto}\textsubscript{G}}, più precisamente nella sezione \textit{Analisi dei rischi} sono contenuti i rischi che potrebbero emergere durante lo svolgimento del progetto, con relativi approfondimenti e strategie di mitigazione.

\subsubsubsection{Codifica dei rischi}
Ogni rischio è identificato da un codice univoco nel seguente formato:
\begin{center}
	\textbf{R[tipologia]-[indice]}: nome identificativo del rischio.
\end{center}
Dove:
\begin{itemize}
	\item \textbf{[tipologia]}: rappresenta la categoria di rischio, la quale può essere organizzativa, tecnologica o comunicativa;
	\item \textbf{[indice]}: un valore numerico incrementale che identifica univocamente il rischio per ogni tipologia.
\end{itemize}

\subsubsection{Metriche}
\begin{table}[h]
	\centering
	\begin{tabular}{|c|c|}
		\hline
		\textbf{Metrica} & \textbf{Nome}        						\\
		\hline
		\underline{\hyperlink{31M}{31M-RMR}}   &  Risk Mitigation Rate  \\
		\underline{\hyperlink{32M}{32M-NCR}}   &  Rischi Non Calcolati  \\
		\hline
	\end{tabular}
	\caption{Metriche relative alla Risoluzione dei problemi}
\end{table}

\subsubsection{Strumenti}
\begin{itemize}
	\item \href{https://7last.github.io/docs/rtb/documentazione-interna/glossario\#clickup}{\textbf{ClickUp}\textsubscript{G}}: utilizzato per la gestione delle attività e delle issue;
\end{itemize}


\newpage


\section{Processi organizzativi}
Lo sviluppo software è un processo complesso e multidisciplinare che richiede una pianificazione, una gestione del tempo e delle risorse accurata, efficiente ed efficace. L'adozione di processi organizzativi ben strutturati è punto cruciale per garantire il successo dello sviluppo software. 
\subsection{Gestione dei processi}
\subsubsection{Introduzione}
La gestione dei processi si occupa di determinare, migliorare, ottimizzare i processi che fanno da guida alla realizzazione del software. Le attività di gestione dei processi sono:
\begin{itemize}
    \item \textbf{Definizione dei processi}:
        \begin{itemize}
            \item Identificare e documentare i processi chiave coinvolti nello sviluppo software
            \item Stabilire le linne guida e procedure per l'esecuzione di ciascun processo
        \end{itemize}
    \item \textbf{Pianificazione e monitoraggio}:
        \begin{itemize}
            \item Elaborare piani dettagliati per l'esecuzione dei processi
            \item Monitorare costantemente l'avanzamento, l'efficacia e la conformità ai requisiti pianificati
            \item Stimare i tempi, le risorse ed i costi
        \end{itemize}
    \item \textbf{Valutazione e miglioramento continuo}
        \begin{itemize}
            \item Condurre valutazioni periodiche dei processi per identificare aree di miglioramento
            \item Implementare azioni correttive e preventive per ottimizzare i processi
        \end{itemize}
    \item \textbf{Formazione e Competenze}
        \begin{itemize}
            \item Assicurare che il personale coinvolto nei processi sia adeguatamente formato
            \item Mantenere e sviluppare le competenze necessarie per l'efficace gestione dei processi
        \end{itemize}
    \item \textbf{Gestione dei rischi}
        \begin{itemize}
            \item Identificare e valutare i rischi associati ai processi
            \item Definire le strategie per mitigare o gestire i rischi identificati
        \end{itemize}
\end{itemize}
\subsubsection{Pianificazione}
\subsubsubsection{Descrizione}
La pianificazione riveste un ruolo centrale nella gestione dei processi, poiché mira a creare un piano organizzato e coerente per assicurare un’efficace esecuzione delle attività durante l’intero ciclo di vita del software.
Il responsabile del progetto assume il compito di coordinare ogni aspetto della pianificazione delle attività, che include l’allocazione delle risorse, la definizione dei tempi e la redazione di piani dettagliati. Inoltre, il responsabile si assicura che il piano elaborato sia fattibile e possa essere eseguito correttamente ed efficientemente dai membri del team.
I piani associati all’esecuzione del processo devono comprendere descrizioni dettagliate delle attività e delle risorse necessarie, specificando le tempistiche, le tecnologie impiegate, le infrastrutture coinvolte e il personale assegnato.
\subsubsubsection{Obiettivi}
L’obiettivo primario della pianificazione è assicurare che ciascun membro del team assuma ogni ruolo almeno una volta durante lo svolgimento del progetto, promuovendo così una distribuzione equa delle responsabilità e un arricchimento delle competenze all’interno del team.
La pianificazione, stilata dal responsabile, è integrata nel documento del \href{https://7last.github.io/docs/rtb/glossario#piano di progetto}{Piano di Progetto\textsubscript{G}}. Questo documento fornisce una descrizione completa delle attività e dei compiti necessari per raggiungere gli obiettivi prefissati in ogni periodo del progetto.

\subsubsubsection{Assegnazione dei ruoli}
Durante l’intero periodo del progetto, i membri del gruppo assumeranno sei ruoli distinti, ovvero assumeranno le responsabilità e svolgeranno le mansioni tipiche dei professionisti nel campo dello sviluppo software.
Nei successivi paragrafi sono descritti in dettaglio i seguenti ruoli:
\begin{itemize}
    \item Responsabile
    \item Amministratore
    \item Analista
    \item Progettista
    \item Programmatore
    \item Verificatore
\end{itemize}
\subsubsubsection{Responsabile}
Figura fondamentale che coordina il gruppo, fungendo da punto di riferimento per il committenteG e il team, svolgendo il ruolo di mediatore tra le due parti.
In particolare si occupa di:
\begin{itemize}
    \item Gestire le relazione con l'esterno
    \item Pianificare le attività: quali svolgere, data di inizio e fine, assegnazione delle priorità
    \item Valutare i rischi delle scelte da effettuare
    \item Controllare i progressi del progetto
    \item Gestire le risorse umane
    \item Approvazione della documentazione
\end{itemize}
\subsubsubsection{Amministratore}
Questa figura professionale è incaricata del controllo e dell’amministrazione dell’ambiente di lavoro utilizzato dal gruppo ed è anche il punto di riferimento per quanto concerne le norme di progetto. Le sue mansioni principali sono:
\begin{itemize}
    \item Affrontare e risolvere le problematiche associate alla gestione dei processi
    \item Gestire versionamento della documentazione 
    \item Gestire la configurazione del prodotto
    \item Redigere ed attuare le norme e le procedure per la gestione della qualità
    \item Amministrare le infrastrutture e i servizi per i processi di supporto
\end{itemize}
\subsubsubsection{Analista}
Figura professionale con competenze avanzate riguardo l’attività di analisi dei requisitiG ed il dominio applicativo del problema. Il suo ruolo è quello di identificare, documentare e comprendere a fondo le esigenze e le specifiche del progetto, traducendole in requisiti chiari e dettagliati. Si occupa di:
\begin{itemize}
    \item Analizzare il contesto di riferimento, definire il problema in esame e stabilire gli obiettivi da raggiungere
    \item Comprendere il  problema e definire la complessità e i requisiti
    \item Redigere il documento \textit{Analisi dei requisiti}
    \item Studiare i bisogni espliciti ed impliciti
\end{itemize}
\subsubsubsection{Progettista}
Il progettista è la figura di riferimento per quanto riguarda le scelte progettuali partendo dal lavoro dell’analista. Spetta al progettista assumere decisioni di natura tecnica e tecnologica, oltre a supervisionare il processo di sviluppo. Tuttavia, non è responsabile della manutenzione del prodotto.
In particolare si occupa di:
\begin{itemize}
    \item Progettare l’architetturaG del prodotto secondo specifiche tecniche dettagliate
    \item Prendere decisioni per sviluppare soluzioni che soddisfino i criteri di affidabilità,
    efficienza, sostenibilità e conformità ai requisiti;
    \item Redige la Specifica Architetturale e la parte pragmatica del \href{https://7last.github.io/docs/rtb/glossario#piano di qualifica}{Piano di Qualifica\textsubscript{G}};
\end{itemize}
\subsubsubsection{Programmatore}
Il programmatore è la figura professionale incaricata della scrittura del codice software. Il suo compito primario è implementare il codice conformemente alle specifiche fornite dall’analista e all’architetturaG definita dal progettista.
In particolare, il programmatore:
\begin{itemize}
    \item Scrive codice manutenibile in conformità con le Specifiche Tecniche
    \item Codifica le varie componenti dell’architettura seguendo quanto ideato dai progettisti
    \item Realizza gli strumenti per verificare e validare il codice
    \item Redige il \textit{Manuale Utente}
\end{itemize}
\subsubsubsection{Verificatore}
La principale responsabilità del verificatore consiste nell’ispezionare il lavoro svolto da altri membri del team per assicurare la qualità e la conformità alle attese prefissate. Stabilisce se il lavoro è stato svolto correttamente sulla base delle proprie competenze tecniche, esperienza e conoscenza delle norme.
In particolare il verificatore si occupa di:
\begin{itemize}
    \item Verificare che il lavoro svolto sia conforme alle Norme di progetto
    \item Verificare che il lavoro svolto sia conforme alle Specifiche Tecniche
    \item Ricercare ed in caso segnalare eventuali errori
    \item Redigere la sezione retrospettiva del \textit{Piano di Qualifica}, descrivendo le verifiche e le prove effettuate durante il processo di sviluppo del prodotto
\end{itemize}
\subsubsubsection{Ticketing}
GitHub è adottato come sistema di tracciamento delle issue (ITS), garantendo così una gestione agevole e trasparente delle attività da svolgere.
L’amministratore ha la facoltà di creare e assegnare specifiche issue sulla base delle attività identificate dal responsabile, assicurando chiarezza sulle responsabilità di ciascun membro del team e stabilendo tempi definiti entro cui ciascuna attività deve essere completata. Inoltre, ogni membro del gruppo può monitorare i progressi compiuti nel periodo corrente, consultando lo stato di avanzamento delle varie issue attraverso le Dashboard:
\begin{itemize}
    \item DashBoard: per una panoramica dettagliata sullo stato delle issue
    \item RoadMap: per una panoramica temporale dettagliata delle issue
\end{itemize}
Procedura per la creazione delle issue:\\
Le issue vengono create dall’amministratore e devono essere specificati i seguenti attributi:
\begin{itemize}
    \item Titolo: breve descrizione dell’attività da svolgere
    \item Descrizione:
        \begin{itemize}
            \item Descrizione testuale oppure "to-do" tramite bullet points
            \item Nell'ultima riga viena specificato il verificatore della issue nel formato: "Verificatore: Nome Cognome"
        \end{itemize}
    \item Assegnatario: membro del team responsabile dell'issue
    \item Milestone: periodo di riferimento in cui l’attività deve essere completata
    \item Labels: etichette per categorizzare le issue. \\Per associare ad ogni issue un Configuration Item vengono utilizzati i seguenti label:
        \begin{itemize}
            \item \textbf{NdP}: Norme di Progetto
            \item \textbf{PdP}: Piano di Progetto
            \item \textbf{PdQ}: Piano di Qualifica
            \item \textbf{AdR}: Analisi dei Requisiti
            \item \textbf{PoC}: Proof of Concept
            \item \textbf{Gls}: Glossario
        \end{itemize}
    \item Milestone: milestone associata alla issue
    \item Projects: progetti a cui la issue è associata. Se sono presenti dashboard associate ad un progetto, le issue correlate a tale progetto verranno visualizzate nella relativa/e dashboard di progetto.
    \item Development: branch e Pull Request associate alla issue. Quando una Pull Request viene accettata, la relativa issue viene automaticamente chiusa ed eventualmente spostata nella sezione "Done" della dashboard di progetto.
\end{itemize}
Ciclo di vita di una issue:\\
Il ciclo di vita è il seguente:
\begin{itemize}
    \item Creazione: l’amministratore crea la issue e la assegna al membro del team responsabile
    \item L’amministratore accede alla dashboardG di progetto e sposta la issue dalla colonna "No Status" alla colonna "To Do";    
    \item L’assegnatario apre un branch su GitHub seguendo la denominazione suggerita in "Sincronizzazione e Branching";
    \item Quando la issue viene presa in carico dall’assegnatario, questo accede alla DashBoard e sposta la issue dalla colonna "To Do" alla colonna "In Progress";
    \item Una volta che la issue è considerata terminata, l’assegnatario apre una Pull Request su GitHub seguendo la convenzione descitta in dettaglio nella sezione "Procedura per la creazione di Pull Request".
    \item All’interno della Dashboard GitHub la issue deve essere spostata dalla colonna "In Progress" alla colonna "Da revisionare";
    \item il verificatore o i verificatori designati seguono le procedure esposte nella sezione 3.2 per verificare le modifiche apportate al progetto;
    \item Se la verifica ha esito positivo, la issueG viene trasferita dalla colonna "Da revisionare" alla colonna "Done" della Dashboard di GitHub. Nel caso in cui la issue sia associata ad una Pull Request, una volta che quest’ultima viene accettata dal verificatore, la issue viene automaticamente chiusa e spostata nella colonna "Done" della Dashboard di progetto.
\end{itemize}
\subsubsubsection{Strumenti}
\begin{itemize}
    \item GitHub: piattaforma utilizzata per il tracciamento e la gestione delle issue.
\end{itemize}
\subsubsection{Coordinamento}
\subsubsubsection{Descrizione}
Il coordinamento rappresenta l’attività che sovraintende la gestione della comunicazione e la pianificazione degli incontri tra le diverse parti coinvolte in un progetto di ingegneria del software. Questo comprende sia la gestione della comunicazione interna tra i membri del team del progetto, sia la comunicazione esterna con il proponente e i committenti. Il coordinamento risulta essere cruciale per assicurare che il progetto proceda in modo efficiente e che tutte le parti coinvolte siano informate e partecipino attivamente in ogni fase del progetto.
\subsubsubsection{Obiettivi}
Il coordinamento in un progetto è fondamentale per gestire la comunicazione e pianificare gli incontri tra gli stakeholder. L’obiettivo principale è garantire efficienza, evitando ritardi e confusioni, assicurando che tutte le parti in causa siano informate e coinvolte in ogni fase del progetto. Inoltre, promuove la collaborazione e la coesione nel team, facilitando lo scambio di idee e la risoluzione dei problemi in modo collaborativo, creando un ambiente lavorativo positivo e produttivo.
\textbf{Comunicazione}
Il gruppo 7Last mantiene comunicazioni attive, sia interne che esterne al team, le quali possono essere sincrone o asincrone, a seconda delle necessità.
\subsubsubsection{Comunicazioni sincrone}
\begin{itemize}
    \item \textbf{Comunicazione sincrone interne}\\Per le comunicazioni sincrone interne, il gruppo 7Last, ha scelto di adottare Discord in quanto permette di comunicare tramite chiamate vocali, videochiamate, messaggi di testo, media e file in chat private o come membri di un "server Discord";
    \item \textbf{Comunicazione sincrone esterne}\\Per le comunicazioni sincrone esterne,in accordo con l’azienda proponente si è deciso di utilizzare un canale Discord.
\end{itemize}
\subsubsubsection{Comunicazioni asincrone}
\begin{itemize}
    \item \textbf{Comunicazione asincrone interne}\\Per le comunicazioni asincrone interne, il gruppo 7Last, ha scelto di adottare Telegram in quanto permette di comunicare tramite messaggi di testo, media e file in chat private o come membri di un "gruppo Telegram";
    \item \textbf{Comunicazione asincrone esterne}\\Per le comunicazioni asincrone esterne sono stati adottati due canali differenti:
        \begin{itemize}
            \item \textbf{Email}: per comunicazioni formali e ufficiali
            \item \textbf{Discord}
        \end{itemize}
\end{itemize}
\subsubsubsection{Riunioni interne}
Si è scelto di svolgere i meeting interni a cadenza settimanale, al fine di facilitare una comunicazione costante e coordinare il progresso delle attività.
Generalmente le riunioni sono programmate per ogni venerdi alle ore:
\begin{itemize}
    \item Mercoledì dalle 15:00 - 16:00 per riunioni interne
\end{itemize}
Se qualche membro del gruppo non può partecipare alla riunione nella data e nell’orario stabiliti, si procede programmando un nuovo incontro, concordando data e ora tramite un sondaggio sul canale Telegram dedicato.
Ogni membro del gruppo ha la facoltà di richiedere una riunione supplementare se necessario. In questo caso, la data e l’orario saranno concordati sempre attraverso il canale Telegram dedicato, mediante la creazione di un sondaggio.
Le riunioni interne rivestono un ruolo cruciale nel monitorare il progresso delle mansioni assegnate, valutare i risultati conseguiti e affrontare i dubbi e le difficoltà che possono sorgere. Durante i meeting interni, i membri del team condividono gli aggiornamenti sulle proprie attività, identificano le problematiche riscontrate e discutono di opportunità di miglioramento nei processi di lavoro. Questo ambiente aperto e collaborativo favorisce l’interazione, l’innovazione e la condivisione di nuove prospettive. Per agevolare la comunicazione sincrona, il canale utilizzato per i meeting interni è Discord, ritenuto particolarmente efficace per tali scopi.
Relativamente ai meeting interni, sarà compito del responsabile:
\begin{itemize}
    \item Stabilire preventivamente i principali temi da trattare durante la riunione, considerando la possibilità di aggiungerne di nuovi nel corso della riunione stessa;
    \item Guidare la discussione e raccogliere i pareri dei membri in maniera ordinata;
    \item Nominare un segretario per la riunione;
    \item Pianificare e proporre le nuove attività da svolgere.
\end{itemize}
\textbf{Verbali interni}\\
Lo svolgimento di una riunione interna ha come obiettivo la retrospettiva del periodo precedente, la discussione dei punti stilati nell’ordine del giorno e la pianificazione delle nuove attività.
Alla conclusione di ciascuna riunione, l’amministratore apre un’issue nell’ITS di GitHub e assegna l’incarico di redigere il verbale interno al segretario della riunione. È compito quindi del segretario redigere il verbale, includendo tutte le informazioni rilevanti emerse durante la riunione. Le indicazioni dettagliate per la compilazione dei verbali interni sono disponibili nella sezione 3.1.6.5.

\subsubsubsection{Riunioni esterne}
Durante lo svolgimento del progetto, è essenziale organizzare vari incontri con i Committenti e/o il Proponente al fine di valutare lo stato di avanzamento del prodotto e chiarire eventuali dubbi o questioni.
La responsabilità di convocare tali incontri ricade sul responsabile, il quale è incaricato di pianificarli e agevolarne lo svolgimento in modo efficiente ed efficace.
Sarà compito del responsabile anche l’esposizione dei punti di discussione al proponente/committente, lasciando la parola ai membri del gruppo interessati quando necessario. Questo approccio assicura una comunicazione efficace tra le varie parti in causa, garantendo una gestione ottimale del tempo e una registrazione accurata delle informazioni rilevanti emerse durante gli incontri.
I membri del gruppo si impegnano a garantire la propria presenza in modo costante alle riunioni, facendo il possibile per riorganizzare eventuali altri impegni al fine di partecipare. Nel caso in cui gli obblighi inderogabili di un membro del gruppo rendessero impossibile la partecipazione, il responsabile assicurerà di informare tempestivamente il proponente o i committenti, richiedendo la possibilità di rinviare la riunione ad una data successiva.

\textbf{Riunioni con l'azienda proponente}\\
In accordo con l’azienda proponente, si è stabilito di tenere incontri di stato avanzamento lavori (SAL) con cadenza bisettimanale tramite Google Meet.
Durante tali incontri, si affrontano diversi aspetti, tra cui:
\begin{itemize}
    \item Discussione delle attività svolte nel periodo precedente, valutando l’aderenza a quanto concordato e identificando eventuali problematiche riscontrate
    \item Pianificazione delle attività per il prossimo periodo, definendo gli obiettivi e le azioni necessarie per il loro raggiungimento
    \item Chiarezza e risoluzione di eventuali dubbi emersi nel corso delle attività svolte
\end{itemize}

\textbf{Verbali esterni}\\
Come nel caso delle riunioni interne, anche per le riunioni esterne verrà redatto un verbale con le stesse modalità descritte nella sezione relativa ai Verbali Interni.
Le linee guida per la redazione dei verbali esterni sono reperibili alla sezione 3.1.6.5.


\subsubsubsection{Strumenti}
\begin{itemize}
    \item \textbf{Discord}: impiegato per la comunicazione sincrona e i meeting interni del team e per le riunioni esterne con il proponente
    \item \textbf{Telegram}: utilizzato per la comunicazione asincrona interna
    \item \textbf{Gmail}: come servizio di posta elettronica
\end{itemize}

\subsubsubsection{Metriche}
AGGIUNGERE TABELLA

\subsection{Miglioramento}
\subsubsection{Introduzione}
Secondo lo standard ISO/IEC 12207:1995, il processo di miglioramento nel ciclo di vita del software è finalizzato a stabilire, misurare, controllare e migliorare i processi che lo compongono. L’attività di miglioramento è composta da:
\begin{itemize}
    \item \textbf{Analisi}: valutazione dei processi per identificare le aree di miglioramento
    \item \textbf{Miglioramento}: attuazione di azioni correttive e preventive per ottimizzare i processi
\end{itemize}
\subsubsection{Analisi}
Questa operazione richiede di essere eseguita regolarmente e ad intervalli di tempo appropriati e costanti. L’analisi fornisce un ritorno sulla reale efficacia e correttezza dei processi implementati, permettendo di identificare prontamente quelli che necessitano di miglioramenti.
Durante ogni riunione, il team dedica inizialmente del tempo per condurre una retrospettiva sulle attività svolte nell’ultimo periodo. Questa pratica implica una riflessione approfondita su ciò che è stato realizzato, coinvolgendo tutti i membri nella identificazione delle aree di successo e di possibili miglioramenti.
L’obiettivo principale è formulare azioni correttive da implementare nel prossimo sprint, promuovendo così un costante feedback e un adattamento continuo per migliorare le prestazioni complessive del team nel corso del tempo.
\subsubsection{Miglioramento}
Il team implementa le azioni correttive stabilite durante la retrospettiva, successivamente valuta la loro efficacia e le sottopone nuovamente a esame durante la retrospettiva successiva.
L’esito di ogni azione correttiva sarà documentato nella sezione "Revisione del periodo precedente" di ogni verbale.

\subsection{Formazione}
\subsubsection{Introduzione}
L’obiettivo di questa iniziativa è stabilire standardG per il processo di apprendimento all’interno del team, assicurando la comprensione adeguata delle conoscenze necessarie per la realizzazione del progetto.
Si prevede che il processo di formazione del Team assicuri che ciascun membro acquisisca una competenza adeguata per utilizzare consapevolmente le tecnologie selezionate dal gruppo per la realizzazione del progetto.

\subsubsection{Metodo di formazione}
\subsubsubsection{Individuale}
Ciascun membro del team si impegnerà in un processo di autoformazione per adempiere alle attività assegnate al proprio ruolo. Durante la rotazione dei ruoli, ogni membro del gruppo condurrà una riunione con il successivo occupante del suo attuale ruolo, trasmettendo le conoscenze necessarie. Al contempo, terrà una riunione con chi ha precedentemente svolto il ruolo che esso assumerà, con l’obiettivo di apprendere le competenze richieste.

\subsubsubsection{Gruppo}
Sono programmate sessioni formative, condotte dalla proponente, al fine di trasferire competenze relative alle tecnologie impiegate nel contesto del progetto. La partecipazione del team a tali riunioni è obbligatoria.




\newpage


\input{standard_qualità.tex}

\newpage


\input{metriche_qualità.tex}

\newpage


%----------------TABLE EXAMPLE----------------%
% \begin{table}[h!]
% 	\centering
% 	\begin{tabular}{|l|l|l|l|} 
% 	 \hline
% 	 Metrica & Nome & Valore ammissibile & Valore ottimo \\  
% 	 \hline
% 	 MPC01 & Earned Value (EV) & $\geq 0$ & $\leq$ EAC \\
% 	 \hline
% 	 MPC02 & Planned Value (PV) & $\geq 0$ & $\leq$ Budget At Completion (BAC) \\ 
% 	 \hline
% 	 MPC03 & Actual Cost (AC) & $\geq 0$ & $\leq$ EAC \\ 
% 	 \hline
% 	 MPC04 & Cost Variance (CV) & $\geq -7.5\%$ & $\geq 0\%$ \\ 
% 	 \hline
% 	 MPC05 & Schedule Variance (SV) & $\geq -7.5\%$ & $\geq 0\%$ \\ 
% 	 \hline
% 	 MPC06 & Estimated at Completion (EAC) & Errore del $\pm 3\%$ rispetto al BAC  & Equivalente al BAC \\ 
% 	 \hline
% 	 MPC07 & Estimate to Complete (ETC)  & $\geq 0$ & $\leq$ EAC \\ 
% 	 \hline
% 	\end{tabular}
%     \caption{Valori delle metriche inerenti al processo di Fornitura}
% 	\label{table:1}
% \end{table}
%---------------------------------------------%


%----------------FIGURE EXAMPLE---------------%

% \begin{figure*}
%     \centering
%     \begin{tikzpicture}
%         \begin{axis}[
%             width=\textwidth,
%             height=0.5\textwidth,
%             ybar,
%             enlargelimits=0,
%             legend style={at={(0.5,1.15)},
%               anchor=north,legend columns=-1},
%             ylabel={Tempo dedicato},    
%             symbolic x coords={Antonio, Davide, Elena, Leonardo, Matteo, Raul, Valerio},
%             xtick=data,
%             % nodes near coords,
%             nodes near coords align={vertical},
%             ]
%         \addplot coordinates {(Antonio,4) (Davide,4) (Elena,4) (Leonardo,4) (Matteo,4) (Raul,4) (Valerio,4)};
%         \addplot coordinates {(Antonio,4) (Davide,4) (Elena,4) (Leonardo,4) (Matteo,4) (Raul,4) (Valerio,4)};
%         \addplot coordinates {(Antonio,4) (Davide,4) (Elena,4) (Leonardo,4) (Matteo,4) (Raul,4) (Valerio,4)};
%         \addplot coordinates {(Antonio,4) (Davide,4) (Elena,4) (Leonardo,4) (Matteo,4) (Raul,4) (Valerio,4)};
%         \addplot coordinates {(Antonio,4) (Davide,4) (Elena,4) (Leonardo,4) (Matteo,4) (Raul,4) (Valerio,4)};
%         \legend{Amministratore, Responsabile, Verificatore, Analista, Progettista}
%         \end{axis}
%     \end{tikzpicture}
% \caption{impegno preventivo per ruolo di ciascun membro del team durante il primo periodo}
% \end{figure*}

%---------------------------------------------%

\end{document}
